% Options for packages loaded elsewhere
\PassOptionsToPackage{unicode}{hyperref}
\PassOptionsToPackage{hyphens}{url}
%
\documentclass[
  11pt,
]{article}
\usepackage{amsmath,amssymb}
\usepackage{iftex}
\ifPDFTeX
  \usepackage[T1]{fontenc}
  \usepackage[utf8]{inputenc}
  \usepackage{textcomp} % provide euro and other symbols
\else % if luatex or xetex
  \usepackage{unicode-math} % this also loads fontspec
  \defaultfontfeatures{Scale=MatchLowercase}
  \defaultfontfeatures[\rmfamily]{Ligatures=TeX,Scale=1}
\fi
\usepackage{lmodern}
\ifPDFTeX\else
  % xetex/luatex font selection
\fi
% Use upquote if available, for straight quotes in verbatim environments
\IfFileExists{upquote.sty}{\usepackage{upquote}}{}
\IfFileExists{microtype.sty}{% use microtype if available
  \usepackage[]{microtype}
  \UseMicrotypeSet[protrusion]{basicmath} % disable protrusion for tt fonts
}{}
\makeatletter
\@ifundefined{KOMAClassName}{% if non-KOMA class
  \IfFileExists{parskip.sty}{%
    \usepackage{parskip}
  }{% else
    \setlength{\parindent}{0pt}
    \setlength{\parskip}{6pt plus 2pt minus 1pt}}
}{% if KOMA class
  \KOMAoptions{parskip=half}}
\makeatother
\usepackage{xcolor}
\usepackage[margin=1in]{geometry}
\usepackage{longtable,booktabs,array}
\usepackage{calc} % for calculating minipage widths
% Correct order of tables after \paragraph or \subparagraph
\usepackage{etoolbox}
\makeatletter
\patchcmd\longtable{\par}{\if@noskipsec\mbox{}\fi\par}{}{}
\makeatother
% Allow footnotes in longtable head/foot
\IfFileExists{footnotehyper.sty}{\usepackage{footnotehyper}}{\usepackage{footnote}}
\makesavenoteenv{longtable}
\usepackage{graphicx}
\makeatletter
\def\maxwidth{\ifdim\Gin@nat@width>\linewidth\linewidth\else\Gin@nat@width\fi}
\def\maxheight{\ifdim\Gin@nat@height>\textheight\textheight\else\Gin@nat@height\fi}
\makeatother
% Scale images if necessary, so that they will not overflow the page
% margins by default, and it is still possible to overwrite the defaults
% using explicit options in \includegraphics[width, height, ...]{}
\setkeys{Gin}{width=\maxwidth,height=\maxheight,keepaspectratio}
% Set default figure placement to htbp
\makeatletter
\def\fps@figure{htbp}
\makeatother
\setlength{\emergencystretch}{3em} % prevent overfull lines
\providecommand{\tightlist}{%
  \setlength{\itemsep}{0pt}\setlength{\parskip}{0pt}}
\setcounter{secnumdepth}{-\maxdimen} % remove section numbering
\newlength{\cslhangindent}
\setlength{\cslhangindent}{1.5em}
\newlength{\csllabelwidth}
\setlength{\csllabelwidth}{3em}
\newlength{\cslentryspacingunit} % times entry-spacing
\setlength{\cslentryspacingunit}{\parskip}
\newenvironment{CSLReferences}[2] % #1 hanging-ident, #2 entry spacing
 {% don't indent paragraphs
  \setlength{\parindent}{0pt}
  % turn on hanging indent if param 1 is 1
  \ifodd #1
  \let\oldpar\par
  \def\par{\hangindent=\cslhangindent\oldpar}
  \fi
  % set entry spacing
  \setlength{\parskip}{#2\cslentryspacingunit}
 }%
 {}
\usepackage{calc}
\newcommand{\CSLBlock}[1]{#1\hfill\break}
\newcommand{\CSLLeftMargin}[1]{\parbox[t]{\csllabelwidth}{#1}}
\newcommand{\CSLRightInline}[1]{\parbox[t]{\linewidth - \csllabelwidth}{#1}\break}
\newcommand{\CSLIndent}[1]{\hspace{\cslhangindent}#1}
\usepackage{newunicodechar}
\usepackage{graphicx}
\renewcommand{\figurename}{Fig.}
\frenchspacing
\usepackage{needspace}
\usepackage{etoolbox}
\ifcsdef{CSLRightInline}{\AtBeginEnvironment{CSLReferences}{\sloppy}\renewcommand{\CSLRightInline}[1]{\parbox[t]{\linewidth - \csllabelwidth}{\raggedright #1}\break}}{}
\usepackage{lineno}
\linenumbers
\AtBeginEnvironment{longtable}{\nolinenumbers}
\AtEndEnvironment{longtable}{\linenumbers}
\AtBeginEnvironment{figure}{\nolinenumbers}
\AtEndEnvironment{figure}{\linenumbers}
\newunicodechar{​}{\allowbreak{}}
\newunicodechar{±}{\ensuremath{\pm}}
\newunicodechar{×}{\ensuremath{\times}}
\newunicodechar{−}{\ensuremath{-}}
\newunicodechar{→}{\ensuremath{\rightarrow}}
\newunicodechar{≈}{\ensuremath{\approx}}
\newunicodechar{≤}{\ensuremath{\leq}}
\newunicodechar{≥}{\ensuremath{\geq}}
\newunicodechar{≫}{\ensuremath{\gg}}
\newunicodechar{⊂}{\ensuremath{\subset}}
\newunicodechar{∩}{\ensuremath{\cap}}
\newunicodechar{∼}{\ensuremath{\sim}}
\newunicodechar{χ}{\ensuremath{\chi}}
\newunicodechar{μ}{\ensuremath{\mu}}
\newunicodechar{σ}{\ensuremath{\sigma}}
\newunicodechar{Δ}{\ensuremath{\Delta}}
\newunicodechar{α}{\ensuremath{\alpha}}
\newunicodechar{β}{\ensuremath{\beta}}
\newunicodechar{λ}{\ensuremath{\lambda}}
\newunicodechar{ν}{\ensuremath{\nu}}
\newunicodechar{¹}{\textsuperscript{1}}
\newunicodechar{²}{\textsuperscript{2}}
\newunicodechar{³}{\textsuperscript{3}}
\newunicodechar{⁴}{\textsuperscript{4}}
\newunicodechar{⁵}{\textsuperscript{5}}
\newunicodechar{⁶}{\textsuperscript{6}}
\newunicodechar{⁷}{\textsuperscript{7}}
\newunicodechar{⁸}{\textsuperscript{8}}
\newunicodechar{⁹}{\textsuperscript{9}}
\newunicodechar{⁰}{\textsuperscript{0}}
\newunicodechar{⁻}{\textsuperscript{$-$}}
\newunicodechar{₀}{\textsubscript{0}}
\newunicodechar{₁}{\textsubscript{1}}
\newunicodechar{₂}{\textsubscript{2}}
\newunicodechar{₃}{\textsubscript{3}}
\newunicodechar{₄}{\textsubscript{4}}
\newunicodechar{₅}{\textsubscript{5}}
\newunicodechar{₆}{\textsubscript{6}}
\newunicodechar{₇}{\textsubscript{7}}
\newunicodechar{₈}{\textsubscript{8}}
\newunicodechar{₉}{\textsubscript{9}}
\ifLuaTeX
  \usepackage{selnolig}  % disable illegal ligatures
\fi
\IfFileExists{bookmark.sty}{\usepackage{bookmark}}{\usepackage{hyperref}}
\IfFileExists{xurl.sty}{\usepackage{xurl}}{} % add URL line breaks if available
\urlstyle{same}
\hypersetup{
  hidelinks,
  pdfcreator={LaTeX via pandoc}}

\author{}
\date{}

\begin{document}

\begin{center}
{\LARGE\bfseries\setlength{\baselineskip}{1.15\baselineskip}
Recall, not verification, is the bottleneck when frontier LLMs elucidate molecular structures from real spectra\par}
\vspace{1.1em}
{\large \textbf{Ilkham Yabbarov}\textsuperscript{a},  \textbf{Rudra Sondhi}\textsuperscript{a}, \textbf{Rodrigo A. Vargas-Hernández}\textsuperscript{a}\par}
\vspace{0.5em}
{\small \textit{\textsuperscript{a} Department of Chemistry and Chemical Biology, McMaster University, Hamilton, Ontario L8S 4L8, Canada.}\par}
\vspace{0.25em}
{\small\textit{E-mail: }\texttt{ilkhamfy@gmail.com}\par}
\end{center}
\vspace{0.6em}

\begingroup\small\setlength{\leftskip}{2em}\setlength{\rightskip}{2em}
\noindent\rule{\dimexpr\linewidth-4em\relax}{0.4pt}\vspace{-0.4em}

Given the molecular formula together with the infrared band list and
¹H/¹³C shift lists exactly as reported in an open-access paper, how
often does a frontier large language model (here, Claude) recover the
correct molecular \emph{constitution}? We find \textbf{28\%} (top-1,
n=194; 95\% CI 22--35) --- far below the near-100\% implied by curated
demonstrations, and \textbf{15\%} once the benchmark's deliberate 50/50
difficulty balance is reweighted to the composition of the corpus it was
drawn from (§3).

The bottleneck is not the model's judgment but its \emph{proposal}. The
model proposes the true structure for only \textbf{34\%} of compounds;
where it does, forward-verification --- predicting each candidate's ¹³C
spectrum and re-ranking by agreement with the observed one --- selects
it \textbf{89\%} of the time (58/65). \textbf{Recall, not verification,
is the wall.}

We release three things: \textbf{IRexp}, the largest openly
redistributable collection of \emph{experimental} infrared band lists
(121,233 records, a third structure-linked); \textbf{IRSpectra-Bench}, a
blind, mechanically scored benchmark of 194 compounds; and a
training-free \textbf{forward-verification} recipe. The two levers are
distinct. Verification re-ranks a fixed candidate set and lifts top-1
from 28\% to 30\%, while \emph{generating wider} moves recall --- 32\%
to 42\%, carrying top-1 from 23\% to 30\% (§5.2, §5.3); both of those
top-1 steps are directional rather than statistically resolved, and the
recall gain is the better-supported effect. Fine-tuning a small
generator on IRexp moves the wall rather than removing it (§5.6).
Masking the spectra drops top-1 from 23\% to 5\% on the same compounds
(§4.6), and Grok 4.6, Gemini 3.7 Flash and GPT-5.6 Sol all verify better
than they generate (§4.7). All data, predictions and code are released.

\vspace{-0.4em}\noindent\rule{\dimexpr\linewidth-4em\relax}{0.4pt}\endgroup

\hypertarget{introduction}{%
\subsection{1. Introduction}\label{introduction}}

Determining a molecule's structure from its spectra is a central,
time-consuming task in synthetic and analytical chemistry. Machine
learning attacks it with encoder--decoder models trained on paired
corpora --- \emph{Spectro} learns ¹H/¹³C/IR → SELFIES from 6,833
molecules\textsuperscript{1} --- and general-purpose LLMs are now
reported to do it off-the-shelf: a non-peer-reviewed 2026 industrial
white paper found that Claude Opus matched or beat commercial
NMR-prediction software in the forward direction (structure→spectrum,
±0.08 ppm ¹H) and ``recovered all eight simpler structures on every
attempt'' in the inverse.\textsuperscript{2} We treat those numbers as a
motivating claim to test against peer-reviewed benchmarks, not as an
established baseline.

That evaluation is narrow: 15 inverse problems on curated single-ring or
two-fragment molecules, NMR only, seven of them given the
\emph{starting-material structure} as a hint, and ``recovery'' scored
leniently over three runs and three ranked candidates. The question a
chemist actually faces --- \emph{take an arbitrary experimental spectrum
from a paper and recover the structure} --- is left open. It needs a
large, diverse, real benchmark; blind, reproducible scoring; and honest
accounting for the methodological choices that inflate or deflate the
apparent score.

We provide all three. \textbf{IRexp} (§2) is, to our knowledge, the
largest openly \emph{redistributable} collection of experimental IR
\textbf{band lists} (121,233 records; 43,060 structure-linked; 33,201
IR+¹H+¹³C+structure); the superlative is scoped to that object type,
since view-only libraries such as SDBS and commercial libraries hold
more structure-linked \emph{spectra} but cannot be redistributed (§2.1).
On it, \textbf{IRSpectra-Bench} (§3) is, to our knowledge, the first
blind, mechanically scored, complexity-stratified evaluation of
frontier-LLM elucidation on real spectra: measured in depth on Claude,
replicated across three other model families (§4.7), and isolating a
large solver-methodology effect with a within-compound control (§4). Its
ground truth is literature structures resolved deterministically
(OPSIN/RDKit) and checked mechanically against the source articles
(560/560 bands on a seed-fixed sample, §2.3; the \emph{expert-chemist}
review is prepared but not yet run, §7); no LLM curates labels or scores
predictions.

Forward-verification elucidation (§5) turns the model's \emph{strong}
direction (forward prediction) against its \emph{weak} one (inverse
regiochemistry); the loop is prior art (§1.1), the decomposition it
enables is ours, run training-free over \textbf{every} benchmark
compound. The finding is a sharp asymmetry: given a candidate set
containing it, the model \emph{verifies} the correct structure 89\% of
the time, yet \emph{proposes} it for only 34\% of compounds --- recall,
not verification, is the wall (Fig. 1). The gain is bounded (top-1
28\%→30\% on n=194; 23\%→30\% on the 60-compound arm where wide
generation was also tested) and, by our own measurements, cannot exceed
a recall/precision ceiling without sharper verification or 2D-NMR data.

\begin{figure}
\centering
\includegraphics[width=6.3in,height=\textheight]{docs/figures/fig_wall.png}
\caption{All 194 benchmark compounds as a single part-to-whole bar. The
true structure is verified top-1 for 58 (green), recalled but mis-ranked
for 7 (vermilion), and never proposed for 129 (grey) --- \emph{the
wall}, 66\% of the bar. Generation recall is 65/194 = 34\%; conditional
verification precision is 58/65 = 89\%; their product is the 30\%
end-to-end top-1. The two rates have different denominators and are not
differenced.}
\end{figure}

Solver and verifier (Fig. S1) are LLM agents on a consumer subscription
--- no fine-tuning, no API spend --- cheap to \emph{re-run} but not
reproducible: the harness pins no model snapshot, exposes no temperature
or seed, and carries an undisclosed product system prompt, and our own
dispatch instruction for the main rounds was not captured, only the
per-compound data it wrapped. An outside group reproduces this
distributionally, not exactly (§8). Scoring is exactly reproducible:
predictions, ground truth and scorers are released, and every number in
the training-free core (§3--§5.3, §4.6) regenerates from them, except
§5.4's HOSE lookup and GNN, which need an nmrshiftdb2 dump we cannot
redistribute (see \emph{Data and code availability}); that GNN and the
§5.6 generator are trained probes, fenced as complements to the
training-free core.

\hypertarget{related-work}{%
\subsubsection{1.1 Related work}\label{related-work}}

\textbf{Trained spectra→structure models.} The dominant line trains
sequence or graph decoders on spectra:
\emph{Spectro}\textsuperscript{1}, a multitask CNN+transformer for
routine 1D-NMR\textsuperscript{3}, set/graph transformers such as
NMRTrans\textsuperscript{4}, and --- closest to our multimodal setting
--- \textbf{NMIRacle}\textsuperscript{5}, conditioned jointly on IR + ¹H
+ ¹³C. They are accurate \emph{in-distribution} but retrained per
modality and dependent on a labelled spectra→structure corpus. That
resource is scarce rather than absent, and one concurrent effort attacks
it the same way we do: NMRTrans builds NMRSpec, a literature-mined
corpus of experimental ¹H/¹³C spectra\textsuperscript{4}. IRexp is the
complement rather than the competitor --- infrared band lists, released
under permissive licences for redistribution (§2.3) --- and the two
together are the open, experimental, multimodal corpus neither supplies
alone.

\textbf{LLMs as elucidators, and how this work differs.}
LLM-orchestrated tools already plan and execute real syntheses
(ChemCrow\textsuperscript{6}, Coscientist\textsuperscript{7}); such
systems \emph{call} an elucidation routine, and this paper measures what
one returns. General-purpose LLMs have been applied
off-the-shelf\textsuperscript{2}, as have multimodal models over
multi-spectral input (SpectraLLM, SpecMol)\textsuperscript{8,9} and
knowledge-enhanced tree-search reasoning\textsuperscript{10};
\textbf{MolPuzzle}\textsuperscript{11} supplies IR+MS+¹H+¹³C puzzles
with the formula given, and \textbf{IR-Agent}\textsuperscript{12}
emulates expert IR interpretation on experimental spectra with a
multi-agent framework. We claim priority on none of them: IR-Agent
improves \emph{how the model reads a spectrum}, where we ask what limits
the outcome once it has.

Prior benchmarks and trained baselines evaluate on simulated or
software-predicted spectra\textsuperscript{1,5} or on curated
single-instrument libraries --- Alberts et al.~use NIST gas-phase
IR\textsuperscript{13,14} --- where IRSpectra-Bench uses spectra
\textbf{as reported by the authors of the source papers}, across
thousands of laboratories and instruments; the contrast is with both
simulation \emph{and} curated uniformity. That line is also the state of
the art for \emph{trained, formula-conditioned IR} elucidation --- an
IR→structure transformer\textsuperscript{13} now at \textbf{63.8\% top-1
/ 84.0\% top-10 on experimental NIST gas-phase spectra} given the
formula, pretrained on 1,399,806 simulated spectra\textsuperscript{14}.
That headline is on a \textbf{6--13-heavy-atom} subset, the range where
our own accuracy is highest (60.5\% top-1 for ≤15 heavy atoms, §4.1), so
on comparable molecules the training-free LLM and the purpose-trained
transformer are closer than the headline numbers suggest. Size is not
the whole story, though, and we do not claim it is: the same work
reports \textbf{59.94\%} on a 5--35-heavy-atom set (n=5,024), which
spans most of our range and falls only four points. What differs there
is the data, not the size --- a single-instrument gas-phase library
against heterogeneous literature-reported band lists. Neither setting
bounds the other. Reported accuracies also swing with inference method
and scoring harness: GPT-4o is scored at \textbf{1.4\%} on MolPuzzle by
the benchmark's own authors\textsuperscript{11}, at \textbf{27.8\%}
under a plain chain-of-thought harness, and at \textbf{57.8\%} with
knowledge-enhanced tree-search reasoning\textsuperscript{10} --- a
factor of forty for one model on one benchmark. Cross-paper numbers are
therefore not comparable, and we fix and release a single pre-registered
RDKit-InChIKey protocol with bootstrap CIs.

\textbf{Computational NMR for structure validation.} Forward
verification is the LLM analog of a workflow chemists already trust ---
compute the spectrum each candidate \emph{would} give and match it to
experiment: DP4 / DP4+ over GIAO-DFT shifts in
solution\textsuperscript{15,16}, \textbf{NMR crystallography} with GIPAW
shifts in the solid state\textsuperscript{17,18}. We swap the
quantum-chemical predictor for a forward LLM, trading accuracy for zero
setup cost, and inherit the principle that \emph{verification by forward
prediction is easier than inverse generation}. \textbf{Computer-assisted
structure elucidation} (CASE) has run the loop since the 1990s and
overturned published assignments\textsuperscript{19,20}; its generator
is a strict enumerator with exhaustive recall by
construction\textsuperscript{20} --- the axis on which we find the LLM
weak, and the axis existing LLM benchmarks leave unmeasured by reporting
a single aggregate score. The decomposition is stable under every
perturbation we could run --- four Claude models, a second chemical
domain, four verifiers, all recall-bound (§4.4, §4.5, §5.4) --- but is
not tested outside the Claude family (§7).

\textbf{The generate-and-forward-verify loop is not ours.}
\textbf{NMR-Solver}\textsuperscript{21} builds it without an LLM,
retrieving and fragment-recombining candidates and ranking them by
¹H/¹³C shifts forward-predicted with NMRNet (¹³C MAE 1.098 ppm); on
\textasciitilde450 experimental literature spectra with the formula
supplied it reports \textbf{52.89\% top-1}, well above our 28.4\%. §5
asks a different question --- \emph{how much of the remaining error is
generation and how much is verification} --- answered by decomposition
(§5.2) rather than by a better solver. That its predictor is roughly
twice as sharp as the \textasciitilde2 ppm LLM forward-predictor §5.1
identifies as the binding constraint, and delivers roughly twice the
top-1, is what our §5.4 ablation predicts and §5.1's separability
measurement implies. \textbf{NMRAgent}\textsuperscript{22} couples
spectral tools to a knowledge graph and validates on newly isolated
natural products; it is complementary --- a best-effort system, where we
measure an off-the-shelf model.

\begin{center}\rule{0.5\linewidth}{0.5pt}\end{center}

\hypertarget{the-irexp-dataset}{%
\subsection{2. The IRexp dataset}\label{the-irexp-dataset}}

\hypertarget{motivation}{%
\subsubsection{2.1 Motivation}\label{motivation}}

An IRexp record is a \textbf{band list} --- peak positions in cm⁻¹
transcribed by an author into a paper's text --- with co-reported ¹H/¹³C
shift lists and, where resolvable, the 2D structure. \textbf{IRexp
contains no absorbance traces.} That is the published form, and what a
language model consumes, but it is a different object from a digitised
spectrum and the two should not be counted against one another.

Among \emph{digitised spectra}, free downloads run to the NIST
WebBook\textsuperscript{23} (\textasciitilde16k) and the Chemotion ELN
deposit\textsuperscript{24} (\textasciitilde2k); AIST
SDBS\textsuperscript{25} is larger (\textasciitilde54k FT-IR, all
structure-linked) but \emph{view-only} (50 spectra/day, no bulk export),
so \textbf{IRexp does not compete there: SDBS alone holds more
structure-linked spectra than IRexp holds structure-linked band lists.}
Among text-derived \emph{band lists} IRexp is the largest \emph{openly
redistributable} collection by record count, a claim deliberately scoped
to that object type.

\hypertarget{construction}{%
\subsubsection{2.2 Construction}\label{construction}}

Experimental sections report per-compound IR band lists with co-reported
¹H/¹³C NMR in a stable textual convention. Mining it is
mature\textsuperscript{26}; we add a pipeline specialised to the
\emph{IR band list}, the modality those tools passed over.

\begin{itemize}
\tightlist
\item
  \textbf{Discovery.} Open-access literature is harvested in bulk from
  the PMC Open-Access Subset\textsuperscript{27} and the Chemotion FT-IR
  deposit (RADAR4Chem).
\item
  \textbf{Extraction.} A deterministic parser segments experimental text
  per compound and extracts IR wavenumbers and ¹H/¹³C shifts, gated
  against scan-range artefacts and prose false-positives.
\item
  \textbf{Resolution.} IUPAC names go to SMILES with
  OPSIN\textsuperscript{28} and RDKit\textsuperscript{29}
  canonicalisation (InChIKey, SELFIES\textsuperscript{30}), with a
  PubChem\textsuperscript{31} fallback for trivial names; handling
  open-access label conventions and narrative artefacts raised structure
  coverage from 24\% to \textbf{35\%}.
\end{itemize}

\hypertarget{contents-and-licensing}{%
\subsubsection{2.3 Contents and
licensing}\label{contents-and-licensing}}

Table 1 gives what the release holds and where each part came from.

\needspace{14\baselineskip}

\textbf{Table 1. IRexp dataset contents and provenance.}

\begin{longtable}[]{@{}
  >{\raggedright\arraybackslash}p{(\columnwidth - 2\tabcolsep) * \real{0.8021}}
  >{\raggedleft\arraybackslash}p{(\columnwidth - 2\tabcolsep) * \real{0.1979}}@{}}
\toprule\noalign{}
\begin{minipage}[b]{\linewidth}\raggedright
field
\end{minipage} & \begin{minipage}[b]{\linewidth}\raggedleft
value
\end{minipage} \\
\midrule\noalign{}
\endhead
\bottomrule\noalign{}
\endlastfoot
experimental IR band-list records & \textbf{121,233} \\
\ldots co-reporting ¹H and/or ¹³C NMR & 87,075 (72\%) \\
\ldots with a resolved 2D structure & \textbf{43,060 (35.5\%)} \\
\ldots of these, with ¹H and/or ¹³C NMR & \textbf{40,702} \\
\ldots of these, with both ¹H and ¹³C (full quadruples) &
\textbf{33,201} \\
& \\
\emph{by provenance:} author-transcribed from PMC-OA text (CC-BY) &
119,345 \\
\emph{by provenance:} peak-picked from Chemotion ELN deposits (CC-BY-SA)
& 1,888 \\
\end{longtable}

The pools differ: PMC records are author-transcribed (median 9 bands),
Chemotion records peak-picked from deposited spectra (median 39), so
users training on band density should treat them apart.
\texttt{scripts/\hspace{0pt}split\_\hspace{0pt}license\_\hspace{0pt}pools.py}
separates them on \texttt{source\_\hspace{0pt}doi} and stamps each
record's licence.

\textbf{Extraction fidelity is measured, not assumed.} On a seed-fixed
random sample of 60 PMC-sourced records
(\texttt{scripts/\hspace{0pt}audit\_\hspace{0pt}extraction.py\ -\/-n\ 60\ -\/-seed\ 0})
we re-fetched each article and checked every recorded wavenumber against
its text: \textbf{560/560 bands and 60/60 records were confirmed}
(Wilson 95\% CI 99.3--100\% and 94--100\%). This bounds
\emph{transcription} error --- hallucinated, mis-parsed or unit-mangled
values --- below 1\% of bands. It does \textbf{not} measure whether the
parser found every IR string in every paper, a recall question requiring
human reading; that audit is prepared but not yet run (§7).

\textbf{\texttt{irexp\_\hspace{0pt}resolved}} (43,060 records, 100\%
structure-linked) is the benchmark-ready split, \textasciitilde6× the
6,833-molecule set used to train Spectro\textsuperscript{1} (Fig. S2).

\begin{center}\rule{0.5\linewidth}{0.5pt}\end{center}

\hypertarget{benchmark-design-irspectra-bench}{%
\subsection{3. Benchmark design
(IRSpectra-Bench)}\label{benchmark-design-irspectra-bench}}

From \texttt{irexp\_\hspace{0pt}resolved} we draw
\textbf{IRSpectra-Bench}, 194 blind elucidation problems, each giving
the \textbf{molecular formula} (as from HRMS), the \textbf{IR band list}
and the \textbf{¹H and ¹³C shift lists}; no name, SMILES or hint.
Main-round ground truths are \textbf{spectrally validated} by an
automated RDKit check (¹³C peaks vs symmetry-unique carbons, formula
match, SELFIES round-trip) excluding merged or incomplete spectra ---
6/140, leaving 134.
\texttt{scripts/\hspace{0pt}validate\_\hspace{0pt}benchmark.py} runs the
deterministic filter over all three rounds, regenerating every
\texttt{clean\_\hspace{0pt}qids.json} from the released questions and
answers.

The filter does not gate on ¹H, and we report what that leaves. In
\textbf{13 of the 194} retained records the reported ¹H integral exceeds
the reference structure's hydrogen count --- residual solvent, water,
exchangeable protons, or a rotamer mixture. They are printed as a
diagnostic, not excluded: the cohort was fixed in advance, and
re-filtering post hoc on a criterion chosen after seeing results is a
degree of freedom a benchmark should not take. We report the sensitivity
instead: dropping all 13 moves the headline from \textbf{28.4\% to
29.3\%} (53/181), +0.9 points, leaving every conclusion unchanged. The
controlled rounds' 60 compounds are fixed, pre-registered sets, audited
the same way but reported rather than filtered (57/60 pass; §7).

\textbf{Difficulty is a declared property of the compound, not a
post-hoc label.} By RDKit ring analysis a compound is \textbf{simple}
iff it has at most two rings, no fused/spiro/bridgehead system and ≤22
heavy atoms; all others are \textbf{complex} (98 simple / 96 complex),
so the split is exhaustive. The threshold is not load-bearing: sweeping
it from 18 to 26 moves the simple-minus-complex top-1 gap only between
36 and 40 points (39.6 as released;
\texttt{scripts/\hspace{0pt}difficulty\_\hspace{0pt}sensitivity.py}).
This axis is distinct from the continuous size gradient of §4.1 (≤15 /
16--25 / \textgreater25 heavy atoms), and InChIKey de-duplication across
rounds prevents leakage.

\textbf{The 50/50 balance is by design, and the corpus is not balanced.}
The sampler fills each stratum to half the round, so ``98 simple / 96
complex'' is a property of the draw and not an observation about the
literature. Under the same eligibility filters the eligible corpus is
\textbf{17.5\% simple and 82.5\% complex} (5,059 / 23,929 of 28,988;
median 26 heavy atoms against the benchmark's 20), so simple compounds
are enriched \textbf{2.9×}
(\texttt{scripts/\hspace{0pt}corpus\_\hspace{0pt}reweight.py}).
Balancing is the right choice for measuring a gradient --- an unbalanced
draw would put almost no mass in the stratum where the model succeeds
--- but it means the headline is a weighted average under weights we
chose. Reweighted to the corpus, top-1 falls from 28.4\% to
\textbf{15.2\%} (95\% CI 10.7--20.4) and recall from 33.5\% to
\textbf{19.8\%} (14.3--25.7). We report the benchmark figure as the
headline because it is the one the released set reproduces, and the
reweighted figure as the better estimate of what a chemist meets on an
arbitrary paper. The reweighting assumes the one sampler filter it
cannot apply --- recovery of the raw ¹H string, which needs a network
fetch per record --- is independent of difficulty.

\textbf{Scoring is mechanical.} A prediction is \emph{correct} if its
RDKit InChIKey connectivity layer (first 14 characters) matches the
reference: we score \emph{constitution}, so correct constitution with
wrong stereochemistry counts as correct. We report the strict
alternative rather than assert it is immaterial --- the \emph{full},
stereochemistry-sensitive InChIKey gives \textbf{21.1\% top-1 and 25.8\%
recovered} (41/194 and 50/194) against 28.4\% / 33.5\%, 7.3 points
lower, so the constitution figure is an upper bound on
full-stereochemistry accuracy. Constitution is the headline because 1D
¹H/¹³C/IR rarely fixes absolute configuration: only 10.3\% (20/194) of
answers carry a \emph{defined} (assigned R/S) stereocentre, and a model
is penalised there for information the prompt never contained
(\texttt{scripts/\hspace{0pt}score\_\hspace{0pt}main.py\ -\/-stereo}).

\textbf{Metrics.} \textbf{Top-1 (exact constitution)} is the fraction
whose best-ranked candidate matches at the connectivity layer;
\textbf{recovered (top-3)} the fraction whose reference appears among
the up-to-three returned candidates, the lenient ``recovery'' protocol
of ref.\textsuperscript{2} --- their metric, under their name.
\textbf{Generation recall} is the fraction whose reference is in the
candidate pool \emph{before} re-ranking, the ceiling any verifier can
reach; it coincides with recovered (top-3) except in the generate-wide
arm of §5.3, where the pool is larger. \textbf{Conditional-on-recall
precision}, over recall-positive compounds alone, isolates verification
from generation. The forward verifier ranks by a symmetric
\textbf{chamfer distance} between predicted and observed ¹³C peak sets
(lower is better, no equal-count requirement); Morgan(2,
2048)\textsuperscript{32} Tanimoto gives a graded scaffold-family
signal. CIs are bootstrap 95\%; model-vs-model differences use McNemar's
exact test\textsuperscript{33} with Holm correction\textsuperscript{34}.

Solvers are frontier LLM agents (Claude Opus), one sub-agent per batch,
closed-book under a consumer subscription: no tools beyond an RDKit
formula check, no ground truth, verified by grep-auditing transcripts.

Formula adherence is high but imperfect --- 91.3\% (115/126) of
forward-verification candidates match the given formula exactly, 76.6\%
across the full top-3 pool --- and \textbf{not uniform across rounds}:
on top-1 answers it is 95.0\% (38/40) in v3 and 90.0\% (18/20) in the
v2-control, against \textbf{77.6\% (104/134) in the headline main round}
--- a structure of the wrong composition, in violation of a constraint
the solver was explicitly handed, for about one main-round compound in
five. Such answers score as misses, so this inflates nothing, but the
main-round arm is the weakest-constrained
(\texttt{scripts/\hspace{0pt}analyze\_\hspace{0pt}misses.py}).

\begin{center}\rule{0.5\linewidth}{0.5pt}\end{center}

\hypertarget{how-well-do-llms-elucidate-real-structures}{%
\subsection{4. How well do LLMs elucidate real
structures?}\label{how-well-do-llms-elucidate-real-structures}}

\hypertarget{headline-performance}{%
\subsubsection{4.1 Headline performance}\label{headline-performance}}

Solver agents work blind from formula + IR + ¹H + ¹³C and return up to
three ranked candidates, in bounded contexts reset every 2--12 compounds
rather than one long context --- the variable §4.3 shows matters. The
benchmark is all \textbf{194 compounds} (134 spectrally-validated + 60
controlled-round;
\texttt{data/\hspace{0pt}benchmark\_\hspace{0pt}main/\hspace{0pt}raw/\hspace{0pt}}).

\needspace{10\baselineskip}

\textbf{Table 2. Headline elucidation performance on IRSpectra-Bench
(n=194)}, overall and by difficulty stratum, with bootstrap 95\% CIs.

\begin{longtable}[]{@{}
  >{\raggedright\arraybackslash}p{(\columnwidth - 6\tabcolsep) * \real{0.4433}}
  >{\raggedleft\arraybackslash}p{(\columnwidth - 6\tabcolsep) * \real{0.1959}}
  >{\raggedleft\arraybackslash}p{(\columnwidth - 6\tabcolsep) * \real{0.1753}}
  >{\raggedleft\arraybackslash}p{(\columnwidth - 6\tabcolsep) * \real{0.1856}}@{}}
\toprule\noalign{}
\begin{minipage}[b]{\linewidth}\raggedright
metric
\end{minipage} & \begin{minipage}[b]{\linewidth}\raggedleft
overall (n=194)
\end{minipage} & \begin{minipage}[b]{\linewidth}\raggedleft
simple (n=98)
\end{minipage} & \begin{minipage}[b]{\linewidth}\raggedleft
complex (n=96)
\end{minipage} \\
\midrule\noalign{}
\endhead
\bottomrule\noalign{}
\endlastfoot
top-1 exact constitution & \textbf{28.4\%} {[}22--35{]} & 48.0\%
{[}39--57{]} & 8.3\% {[}3--15{]} \\
recovered (within top-3) & 33.5\% {[}27--40{]} & 54.1\% {[}44--63{]} &
12.5\% {[}6--20{]} \\
scaffold-level (best Tanimoto ≥ 0.45) & 56\% & 73\% & 39\% \\
mean best Tanimoto & 0.59 & 0.73 & 0.45 \\
\end{longtable}

A strictly-validated subset with the controlled rounds restricted to
their clean parts too (134 main-clean + 57 controlled-clean = 191)
reproduces this within \textasciitilde1 point on every metric --- top-1
\textbf{28.8\%}, 95\% CI 23--36; recall 34.0\%; simple 49.0\% / complex
8.4\% --- so the asymmetric inclusion of the pre-registered controlled
sets does not drive it. What does move the overall figure is the
deliberate 50/50 difficulty balance: reweighted to the 17.5\%-simple
composition of the eligible corpus, top-1 is \textbf{15.2\%} {[}10.7,
20.4{]} and recall \textbf{19.8\%} {[}14.3, 25.7{]} (§3). The
per-stratum figures are unaffected --- only the weights change --- and
the reweighted number, not the headline, is the one to read as
performance on an arbitrary paper. Accuracy falls monotonically with
size in step with the difficulty gradient (Fig. 2): \textbf{60.5\%}
top-1 at ≤15 heavy atoms, 28.3\% at 16--25, \textbf{7.0\%} above 25
(Fig. S3).

\begin{figure}
\centering
\includegraphics[width=3.3in,height=\textheight]{docs/figures/fig1_difficulty.png}
\caption{Top-1 and recovered accuracy on IRSpectra-Bench by difficulty
(all / simple / complex, n=194), with bootstrap 95\% confidence
intervals. The benchmark separates a realistic difficulty range: simple
targets are solved four to six times as often as complex ones on both
metrics.}
\end{figure}

Of the 137 analysable top-1 misses (139 in all; two predictions did not
parse; \texttt{scripts/\hspace{0pt}analyze\_\hspace{0pt}misses.py}),
\textbf{76.6\% are constitutional isomers of the true structure} against
23.4\% with the wrong formula. That is not mostly regiochemistry: only
\textbf{22.6\%} share the true Murcko scaffold, 2.9\% reach Tanimoto ≥
0.85, and the median isomeric-miss Tanimoto is 0.39 --- strict
regiochemistry is a fifth of failures, not the bulk, most misses being
larger rearrangements of the same atoms. Near-degenerate ¹H/¹³C shifts
under-determine connectivity, hence HMBC and NOESY; the honest statement
is \emph{wrong connectivity at the right composition}.

\hypertarget{reconciling-with-prior-reports}{%
\subsubsection{4.2 Reconciling with prior
reports}\label{reconciling-with-prior-reports}}

Our 28\% top-1 sits far below the \textasciitilde100\% on ``simple''
molecules reported for the same model class in a non-peer-reviewed
company white paper.\textsuperscript{2} Unscored independently, it is
best read against the \emph{peer-reviewed} record:
MolPuzzle\textsuperscript{11} and its re-scorings\textsuperscript{10},
and the trained baselines\textsuperscript{1,5}. Four methodology and
scoring choices differ, each in the direction that raises a reported
number; we name them without apportioning the gap. \textbf{Difficulty:}
ring count bounds difficulty from below rather than proxying for it ---
our simple stratum includes a hexasubstituted benzene, while ≥4 rings
mark 38\% of recall-negative against 3\% of recall-positive compounds
(§5.3). \textbf{Scoring:} prior work counts a recovery if the reference
appears among three ranked candidates over three independent runs; we
report single-run top-1 and top-3. \textbf{Hints:} prior hard targets
got the starting-material SMILES and we give none, though the formula we
do give is worth 5\% top-1 alone with the spectra masked (§4.6).
\textbf{Curation:} prior compounds were hand-selected; ours scraped,
unfiltered for solvability.

\textbf{Versus trained models --- a bound, not a leaderboard.} No system
has been scored on our test set, and published numbers differ in the
three respects that most move the score: spectrum realism, hints,
exact-match definition. Most trained baselines report in-distribution
accuracy on simulated spectra; Alberts et al.~do not --- 63.8\% top-1
from IR alone, formula supplied, on \textbf{experimental} NIST gas-phase
spectra of a curated single-instrument library of
\textbf{6--13-heavy-atom} molecules\textsuperscript{14}. Ours span 8--60
heavy atoms (median 20), only \textbf{15 of 194 (8\%)} in that window;
against our ≤15-heavy-atom stratum (60.5\%) their 63.8\% is a near-tie.
Size alone does not explain the rest: the same work reports
\textbf{59.94\%} over 5--35 heavy atoms (n=5,024), a range that covers
most of ours. The remaining difference is what the spectra are --- one
gas-phase instrument against thousands of laboratories --- and the
comparison bounds the gap between those regimes rather than ranking the
two systems. Multimodal systems report more still, on data of their own
making: Spectro \textbf{93\%} (\textbf{82\%} with fixed embeddings) on a
split whose IR is plotted from reference data and whose NMR is
software-\emph{predicted}\textsuperscript{1}; NMIRacle 48\% top-1 / 66\%
top-15 on molecules held out from a \emph{simulated} corpus inside its
own training distribution, as its authors note\textsuperscript{5}. Ours
is 28.4\% top-1 (33.5\% top-3), 29.9\% with forward verification (§5),
on \textbf{blind, real, literature-mined experimental} spectra of
out-of-distribution compounds.

One asymmetry cuts against us: NMIRacle takes \textbf{no molecular
formula} where we supply it (§3), so the settings differ on opposite
axes --- our spectra real and out-of-distribution against their
simulated and in-distribution, their input strictly harder --- and
neither number bounds the other. Read only as a \emph{bound on the
simulated-to-real gap}, the contrast suggests high in-distribution
accuracies substantially overstate real-world performance. Metric
instability reinforces the caution: the \textasciitilde40× MolPuzzle
swing in §1.1 (GPT-4o: 1.4\%\textsuperscript{11} to 27.8\% to
57.8\%\textsuperscript{10}, method- and harness-dependent) bounds the
weight any unaudited near-100\% claim\textsuperscript{2} can bear.

\hypertarget{methodology-dominates-a-within-compound-control}{%
\subsubsection{4.3 Methodology dominates: a within-compound
control}\label{methodology-dominates-a-within-compound-control}}

The same 20 molecules were solved two ways: (a) one LLM context handling
all sequentially with no tools, (b) four independent agents of five
compounds each with RDKit formula-checking. \textbf{Recovered (top-3)}
rose from \textbf{5\% to 15\%} (1/20 → 3/20) and top-1 from 0\% to 15\%
(0/20 → 3/20). McNemar's exact test (§3) does not reach significance,
and at this size it could not have: the exact test conditions on the
discordant pairs, of which the top-1 arm has three (b=0, c=3,
\textbf{p=0.25}) and the recovered arm four (b=1, c=3,
\textbf{p=0.625}). Below six discordant pairs the smallest attainable
two-sided p is above 0.05, so ``not significant'' here reports the
design, not the data. The 15\% carries a Wilson 95\% CI of roughly
5--36\%. Bounded, frequently-reset contexts with tool access therefore
appear to raise measured performance, consistent in direction but
\textbf{not established in size} at n=20 (p=0.25): a directional
within-compound demonstration, not a measured 3× effect. Small rounds
swing widely (15--40\% across n=20--40 draws), hence the full
\textbf{194-compound} headline. It plausibly explains \emph{part} of the
gap to optimistic prior reports, whose per-problem API calls implicitly
used method (b) --- a hypothesis, not a quantified contribution.

\hypertarget{model-comparison-the-benchmark-ranks-capability-and-separates-the-extremes}{%
\subsubsection{4.4 Model comparison: the benchmark ranks capability (and
separates the
extremes)}\label{model-comparison-the-benchmark-ranks-capability-and-separates-the-extremes}}

We solved a fixed 24-compound subset blind with four Claude models
spanning a wide capability range, including the newest (Fable 5), under
one prompt, one scorer and one candidate budget (Fig. 3, Table 3).

One protocol asymmetry must be disclosed: the three comparison models
solved the 24 compounds as four six-compound contexts each, whereas the
Opus column reuses the headline run, where those items sat in one
six-compound and two twelve-compound contexts --- the variable §4.3
measures at 5\%→15\%. The Opus estimate is therefore not
protocol-matched and, if anything, handicapped by the longer contexts;
this touches neither the nesting nor the Fable-vs-Haiku contrast, and a
clean re-run is outstanding in \texttt{docs/\hspace{0pt}MODELS.md}.

\begin{figure}
\centering
\includegraphics[width=3.3in,height=\textheight]{docs/figures/fig5_models.png}
\caption{Four-model comparison on a fixed 24-compound subset, solved
blind under one protocol: Fable 5 46\% \textgreater{} Opus 25\%
\textgreater{} Sonnet 21\% \textgreater{} Haiku 0\% top-1. The outcomes
are strictly nested --- each stronger model solves a superset of the
weaker one's compounds --- so the benchmark is capability-sensitive, but
at n=24 it is underpowered to separate adjacent models (§4.4).}
\end{figure}

\needspace{10\baselineskip}

\textbf{Table 3. Four-model comparison on a fixed 24-compound subset}
under the identical blind protocol.

\begin{longtable}[]{@{}
  >{\raggedright\arraybackslash}p{(\columnwidth - 6\tabcolsep) * \real{0.3333}}
  >{\raggedleft\arraybackslash}p{(\columnwidth - 6\tabcolsep) * \real{0.1458}}
  >{\raggedleft\arraybackslash}p{(\columnwidth - 6\tabcolsep) * \real{0.2292}}
  >{\raggedleft\arraybackslash}p{(\columnwidth - 6\tabcolsep) * \real{0.2917}}@{}}
\toprule\noalign{}
\begin{minipage}[b]{\linewidth}\raggedright
model
\end{minipage} & \begin{minipage}[b]{\linewidth}\raggedleft
top-1
\end{minipage} & \begin{minipage}[b]{\linewidth}\raggedleft
recovered
\end{minipage} & \begin{minipage}[b]{\linewidth}\raggedleft
top-1 95\% CI
\end{minipage} \\
\midrule\noalign{}
\endhead
\bottomrule\noalign{}
\endlastfoot
Claude Fable 5 & \textbf{46\%} & 54\% & {[}25, 67{]} \\
Claude Opus & 25\% & 29\% & {[}8, 42{]} \\
Claude Sonnet & 21\% & 25\% & {[}8, 38{]} \\
Claude Haiku & \textbf{0\%} & 4\% & {[}0, 14{]} \\
\end{longtable}

CIs are bootstrap except Haiku's: Clopper--Pearson exact for 0/24, where
the percentile bootstrap is degenerate.

The models rank in \textbf{monotonic capability order} with
\textbf{strictly nested} outcomes (Haiku ⊂ Sonnet ⊂ Opus ⊂ Fable), so
the benchmark is capability-sensitive. The nesting makes the
\emph{ranking} robust, but at n=24 the subset is \textbf{underpowered to
separate adjacent models}: by McNemar's exact test only Fable-vs-Haiku
survives multiple-comparison correction (Holm-adjusted p=0.006);
Fable-vs-Opus does not (uncorrected p=0.063 --- which is the
\emph{floor} for its five discordant pairs, so the comparison had no way
to reach 0.05 --- Holm-adjusted p=0.19), and Opus and Sonnet are
indistinguishable. Two mid-tier models agreeing closely puts the
recall-bound regime of §4.1 beyond a single model; the newest nearly
\textbf{doubles} the next-best top-1 yet still misses the majority,
leaving the benchmark far from saturated. We used Claude-family models
because they are callable for free under one subscription; a
cross-vendor sweep needs API access we did without, though the large
monotonic spread makes it unlikely the pattern is lineage-specific.
\#\#\# 4.5 Domain case study: battery-electrolyte chemistry

To test the recall-bound regime on battery-electrolyte chemistry rather
than on organic molecules at large, we curated
\textbf{IRSpectra-Bench-Electrolyte}: a 48-compound subset of IRexp
records (eight per class; two yielded no parseable candidate, leaving
\textbf{46 scored}) across the six functional families that dominate
lithium- and sodium-battery electrolytes, held out of every other split
and solved blind under the identical protocol. These are literature
compounds, \textbf{not} operando or in-cell spectra, and we make no
claim that they are.

Performance lands in the \textbf{same broad regime as the headline
benchmark} --- \textbf{top-1 26\%, recovered 28\%} (12/46 and 13/46; at
n=46 the 95\% CI is wide and overlaps the headline) --- consistent with
a bottleneck in the elucidation task rather than in any one chemical
neighbourhood. The true structure enters the candidate set for only 13
of 46 compounds and ranks first in 12 of those 13. Forward-verification
was \textbf{not} run here, so this is the recall-bound pattern under the
solver's own ranking, not the §5 measurement (Fig. S4). The per-class
breakdown is Table 4; at eight compounds a class the intervals overlap
throughout, so the ordering is hypothesis-generating rather than an
established ranking.

\needspace{12\baselineskip}

\textbf{Table 4. Per-class performance on IRSpectra-Bench-Electrolyte
(n=46).}

\begin{longtable}[]{@{}
  >{\raggedright\arraybackslash}p{(\columnwidth - 10\tabcolsep) * \real{0.2990}}
  >{\raggedleft\arraybackslash}p{(\columnwidth - 10\tabcolsep) * \real{0.0825}}
  >{\raggedleft\arraybackslash}p{(\columnwidth - 10\tabcolsep) * \real{0.0928}}
  >{\raggedleft\arraybackslash}p{(\columnwidth - 10\tabcolsep) * \real{0.1340}}
  >{\raggedleft\arraybackslash}p{(\columnwidth - 10\tabcolsep) * \real{0.2577}}
  >{\raggedleft\arraybackslash}p{(\columnwidth - 10\tabcolsep) * \real{0.1340}}@{}}
\toprule\noalign{}
\begin{minipage}[b]{\linewidth}\raggedright
electrolyte class
\end{minipage} & \begin{minipage}[b]{\linewidth}\raggedleft
n
\end{minipage} & \begin{minipage}[b]{\linewidth}\raggedleft
top-1
\end{minipage} & \begin{minipage}[b]{\linewidth}\raggedleft
95\% CI
\end{minipage} & \begin{minipage}[b]{\linewidth}\raggedleft
recovered (top-3)
\end{minipage} & \begin{minipage}[b]{\linewidth}\raggedleft
95\% CI
\end{minipage} \\
\midrule\noalign{}
\endhead
\bottomrule\noalign{}
\endlastfoot
sp³-C--F & 8 & \textbf{50\%} & {[}22, 78{]} & 50\% & {[}22, 78{]} \\
carbonate & 7 & 29\% & {[}8, 64{]} & 43\% & {[}16, 75{]} \\
phosphoryl & 8 & 25\% & {[}7, 59{]} & 25\% & {[}7, 59{]} \\
glyme / oligoether & 7 & 29\% & {[}8, 64{]} & 29\% & {[}8, 64{]} \\
sulfonyl / sulfonate & 8 & \textbf{12\%} & {[}2, 47{]} & 12\% & {[}2,
47{]} \\
nitrile & 8 & \textbf{12\%} & {[}2, 47{]} & 12\% & {[}2, 47{]} \\
\end{longtable}

The intervals overlap almost completely and the spread is within
sampling noise (six-class χ² p≈0.56; best-vs-worst Fisher exact p≈0.28),
so the ordering is hypothesis-generating rather than an established
ranking, though chemically legible: C--F couplings pin sp³-C--F
regiochemistry, while sulfonyl and nitrile targets turn on
oxidation-state ambiguity and heteroaromatic substitution that ¹H/¹³C
shifts underdetermine. It is the recall-bound failure of §4.1 in a
domain-specific guise, and it is where §5's
\textbf{forward-verification} recipe plays a role \emph{analogous} to
(not a replacement for) computational-NMR validation.

\hypertarget{is-the-model-reading-the-spectra-a-formula-only-control}{%
\subsubsection{4.6 Is the model reading the spectra? A formula-only
control}\label{is-the-model-reading-the-spectra-a-formula-only-control}}

Every benchmark compound is mined from open-access literature, so a
frontier model may have met it in pretraining --- a well-documented
hazard for LLM evaluation.\textsuperscript{35} We reran the identical
blind protocol with every spectral channel masked, leaving the solver
the molecular formula and nothing else, which does not determine
constitution: accuracy materially above the floor would indicate recall,
not reasoning.

The arms were not generated together: the \textbf{formula-only} arm was
fresh (2026-07-28), the \textbf{full-modality} arm archived from June.
The bias therefore runs \emph{toward} formula-only --- fresh agents
reason harder per compound than the archived batched run --- and it
still collapsed to 5\%; a confound that works against the finding cannot
manufacture it.

\needspace{8\baselineskip}

\textbf{Table 5. Formula-only control}, paired on the same 60 compounds
as §5.

\begin{longtable}[]{@{}
  >{\raggedright\arraybackslash}p{(\columnwidth - 4\tabcolsep) * \real{0.4375}}
  >{\raggedleft\arraybackslash}p{(\columnwidth - 4\tabcolsep) * \real{0.2292}}
  >{\raggedleft\arraybackslash}p{(\columnwidth - 4\tabcolsep) * \real{0.3333}}@{}}
\toprule\noalign{}
\begin{minipage}[b]{\linewidth}\raggedright
condition
\end{minipage} & \begin{minipage}[b]{\linewidth}\raggedleft
top-1
\end{minipage} & \begin{minipage}[b]{\linewidth}\raggedleft
recovered (top-3)
\end{minipage} \\
\midrule\noalign{}
\endhead
\bottomrule\noalign{}
\endlastfoot
formula only & \textbf{3/60 (5\%)} & 3/60 (5\%) \\
formula + IR + ¹H + ¹³C & \textbf{14/60 (23\%)} & 19/60 (32\%) \\
\end{longtable}

The outcomes are perfectly nested: \textbf{eleven} compounds are solved
with the spectra and not without, and \textbf{none} the other way round
(McNemar exact p=0.001). Removing the spectra removes the result.

\textbf{A second, independent control: publication recency.} The
formula-only arm shows the spectra carry the signal, not whether
memorisation contributes to the part they explain; if recall drove the
headline number, accuracy should fall with recency. Across \textbf{all
194} compounds, source papers spanning 2008--2026, accuracy is flat:
\textbf{28.6\%} for the older half (≤2020, n=112) against
\textbf{28.0\%} for the newer (n=82), a point-biserial year--correctness
correlation of \textbf{r = −0.007}, and no monotone trend across buckets
(Fig. 4b), the most recent (≥2024, n=25) being highest at 40\% {[}23,
59{]}.

\begin{figure}
\centering
\includegraphics[width=6.3in,height=\textheight]{docs/figures/fig_contamination.png}
\caption{Two contamination controls. (\textbf{a}) The solver reaches
3/60 on formula alone against 14/60 with IR + ¹H + ¹³C, nested outcomes:
11 solved only with the spectra, none only without (McNemar exact
p=0.001). (\textbf{b}) Accuracy against source-paper year (n=194, Wilson
95\% CIs) is flat, point-biserial r = −0.007, where pretraining recall
predicts a decline.}
\end{figure}

Newer papers skew larger (median 22 heavy atoms against 20) and size
dominates accuracy, biasing the raw split \emph{against} them (§4.1);
size-adjusted, the older-minus-newer difference is \textbf{−5.1 points,
95\% CI {[}−17.2, +7.0{]}} (Cochran--Mantel--Haenszel χ²=0.42, p=0.51,
continuity-corrected). Its sign is opposite to what pretraining recall
would produce, but the interval comfortably includes zero, so this
\textbf{bounds} any recency effect rather than demonstrating a reversed
one, and by the standard §5.3 applies to its own adjacent conditions we
do not read it as directional.

The 5\% is not zero, and its three compounds do not all read the same
way: one is \textbf{absent from PubChem} and near-determined by its
benzyne-precursor composition, while a catalogued sulfonamide and a
named natural product are as consistent with recognition as with
inference. Either way, formula-level recall accounts for at most about a
fifth of the headline accuracy on this set. The two controls are
independent and agree, but neither is randomised --- publication year is
observational, and the formula-only arm cannot distinguish memorisation
from inference on a near-determining formula --- so we claim a strong
bound rather than exclusion (§7). Nor can either control address
retrieval on the spectral string itself, which the prompt carries
verbatim from the source article; §7 states that objection and names the
control that would settle it.

\hypertarget{does-the-diagnosis-hold-outside-one-vendor-a-four-vendor-replication}{%
\subsubsection{4.7 Does the diagnosis hold outside one vendor? A
four-vendor
replication}\label{does-the-diagnosis-hold-outside-one-vendor-a-four-vendor-replication}}

Every number so far comes from one lineage, the external-validity gap §7
(iii) named as most important. Six non-Claude models solved the same 60
compounds as Table 7's first column under the identical blind protocol;
the three whose output was well-formed enough were carried through
forward verification (Table 6).

\needspace{10\baselineskip}

\textbf{Table 6. Cross-vendor decomposition on the 60-compound arm.}
Recall and precision have different denominators, so the criterion is
the inequality, not a difference.

\begin{longtable}[]{@{}
  >{\raggedright\arraybackslash}p{(\columnwidth - 6\tabcolsep) * \real{0.2917}}
  >{\raggedleft\arraybackslash}p{(\columnwidth - 6\tabcolsep) * \real{0.2708}}
  >{\raggedleft\arraybackslash}p{(\columnwidth - 6\tabcolsep) * \real{0.2708}}
  >{\raggedleft\arraybackslash}p{(\columnwidth - 6\tabcolsep) * \real{0.1667}}@{}}
\toprule\noalign{}
\begin{minipage}[b]{\linewidth}\raggedright
model
\end{minipage} & \begin{minipage}[b]{\linewidth}\raggedleft
generation recall
\end{minipage} & \begin{minipage}[b]{\linewidth}\raggedleft
verification precision \textbar{} recall
\end{minipage} & \begin{minipage}[b]{\linewidth}\raggedleft
multi-candidate only
\end{minipage} \\
\midrule\noalign{}
\endhead
\bottomrule\noalign{}
\endlastfoot
Claude Opus (Table 7) & 19/60 = 32\% {[}21, 44{]} & 16/19 =
\textbf{84\%} {[}62, 94{]} & 10/13 = 77\% \\
Grok 4.6 & 32/60 = \textbf{53\%} {[}41, 65{]} & 20/32 = 62\% {[}45,
77{]} & 20/32 = 62\% \\
Gemini 3.7 Flash & 30/60 = \textbf{50\%} {[}38, 62{]} & 22/30 = 73\%
{[}56, 86{]} & 22/30 = 73\% \\
GPT-5.6 Sol & 25/60 = \textbf{42\%} {[}30, 54{]} & 17/25 = 68\% {[}48,
83{]} & 16/24 = 67\% \\
\end{longtable}

\textbf{The inequality holds in every arm} (Fig. 5): verification
precision exceeds generation recall for four independent model families.
Recall and precision are measured on the same compounds, so what carries
the claim is the paired difference, not whether two marginal intervals
overlap. Bootstrapping compounds
(\texttt{scripts/\hspace{0pt}cross\_\hspace{0pt}vendor\_\hspace{0pt}gap.py})
resolves three of the four: Claude, Gemini \textbf{+23.3 points {[}+2.9,
+44.3{]}} and GPT-5.6 Sol \textbf{+26.3 {[}+4.2, +49.0{]}} all exclude
zero. Grok does not --- \textbf{+9.2 {[}−10.7, +30.3{]}} --- and we
report its gap as directional.

Two findings sit underneath. Six of Claude's nineteen recall-positive
compounds carried one candidate, so nothing was ranked and any verifier
scores them by construction --- as §5.2 says, but only against itself
--- whereas the newer models almost always return three (0--1
singletons), and where a choice existed the four precisions fall in a
band that dissolves the apparent Claude advantage. And verification
looks vendor-independent where generation does not, precision on
multi-candidate sets spanning 62--77\% while recall spans 32--53\%: the
wall is the generator. Three models beat ours on generation recall,
which is the diagnosis working: recall is the movable factor.

\textbf{The candidate budget is not matched, and the correction runs
both ways.} Recall above asks whether the true structure is anywhere in
the candidate list, and the lists differ in length: ours holds
\textbf{2.20} candidates per compound on this arm against exactly
\textbf{3.00} for every comparison model (Composer 2.82). A longer list
can only raise recall, so the comparison is tilted toward the other
vendors --- the mirror of the singleton effect we correct on the
precision side. At \textbf{one} candidate, the only budget every model
met, recall is 14/60 = 23\% for ours against 23/60 = 38\% for Grok,
23/60 = 38\% for Gemini and 21/60 = 35\% for GPT-5.6 Sol
(\texttt{scripts/\hspace{0pt}cross\_\hspace{0pt}vendor\_\hspace{0pt}budget.py}).
The ordering survives and three models still lead ours, but Grok's
margin is \textbf{15 points, not 21}, and 21 is the figure a
budget-matched comparison does not support.

Three weaker models are excluded from the decomposition: Composer 2.5
(20\% recall) and GPT-5.6 Luna (15\%) match the given formula only 67\%
and 76\% of the time, and
\texttt{nvidia/\hspace{0pt}nemotron-3.5-lightning} managed 2\%. §4.1's
scorer now prints formula adherence beside recall. A seventh arm,
DeepSeek V4 Pro, returned answers for only \textbf{18 of 60} compounds
before the run ended; the true structure appears for 8 of those, so its
recall is a lower bound and it is excluded from the comparison rather
than from the release --- both its solve and verify files are deposited.

\begin{figure}
\centering
\includegraphics[width=6.3in,height=\textheight]{docs/figures/fig7_crossvendor.png}
\caption{Every model measured on the 60-compound arm. (\textbf{a})
Generation recall, each model at the candidate budget it used --- 3.00
candidates per compound for all but ours (2.20) and Composer (2.82), so
this axis favours the models that always returned three; grey marks
models below the formula-adherence floor, whose recall is not a
chemistry result. (\textbf{b}) The claim itself: verification precision
against generation recall, with the diagonal. Every model sits above the
line, so verification is better than generation for all four families;
hollow marks an incomplete arm, where recall is a lower bound. The four
Claude models of §4.4 are deliberately absent --- they ran a different
24-compound subset, and putting them here would imply a comparison that
was never made.}
\end{figure}

\textbf{Contamination was controlled, not assumed.} Blindness rested on
an instruction: these runs used cloud agents with repository access to
tracked answer keys. Grok re-solved all ten batches from a clone with
the answer files removed: recall 28/60 against 32/60, paired McNemar
\textbf{p=0.39}, formula adherence 97\% against 95\% of candidates
(\texttt{scripts/\hspace{0pt}cross\_\hspace{0pt}vendor\_\hspace{0pt}sweep.py}).
The asymmetry is decisive --- a key-reader would solve a
\emph{superset}, yet \textbf{four compounds were solved only in the arm
that had no key}, with 24 of 60 solved by both, which is sampling noise,
not copying. And on constitutionally correct structures whose reference
carries assigned stereocentres, which 1D spectra cannot supply, Grok
reproduced 0/3 correct descriptors, Gemini 0/2 and GPT-5.6 Sol 0/2.

Two limits. The models were reached through a coding-assistant harness
exposing reasoning-effort tiers, a decoding control the Claude arm never
had (§8), so these measure a model \emph{as served}, not a bare
endpoint. And the clean-clone control covers only Grok; Gemini and
GPT-5.6 Sol rest on the stereochemistry argument and the shared
protocol.

\begin{center}\rule{0.5\linewidth}{0.5pt}\end{center}

\hypertarget{forward-verification-elucidation}{%
\subsection{5. Forward-verification
elucidation}\label{forward-verification-elucidation}}

\hypertarget{method}{%
\subsubsection{5.1 Method}\label{method}}

The inverse direction is the model's hard, isomer-blind direction; the
\textbf{forward} direction (structure→spectrum) is its easy, accurate
one.\textsuperscript{2} We close a training-free generator--verifier
loop:

\begin{quote}
\textbf{generate} candidate structures (inverse) →
\textbf{forward-predict} each candidate's ¹³C spectrum, blind to the
observed spectrum → \textbf{re-rank} candidates by the distance between
predicted and observed ¹³C → return the best match.
\end{quote}

Candidates are the inverse solver's proposals: 373 deduplicated
structures across the 194 targets, forward-predicted in shuffled,
anonymised batches, blind to the observed spectrum and the target's
identity. Predicted and observed ¹³C peak sets are then compared by
symmetric chamfer distance (Fig. 6).

\begin{figure}
\centering
\includegraphics[width=6.3in,height=\textheight]{docs/figures/fig_mechanism.png}
\caption{Forward-verification on a benchmark regioisomer pair:
picolinamide and nicotinamide are indistinguishable to the inverse task,
but the true isomer's forward-predicted ¹³C spectrum matches the
observed one at a chamfer of 0.42 ppm against 1.30 ppm for the
alternative (§5.1).}
\end{figure}

Regioisomers have \emph{different} forward-predicted shifts, but the
margin is thin: isomeric candidate pairs for the same target separate by
a median \textbf{1.21 ppm} (quartiles 0.84 and 1.78), and \textbf{82\%
are predicted closer together than the predictor's own \textasciitilde2
ppm error} --- usually within the noise. The ranking nevertheless
recovers the right one 89\% of the time when it is present (§5.2), a
real effect on a thin margin, confirmed against a derangement null at
p=0.001 (§5.5).

Conditional precision is 72--89\%, not 100\%: near-degenerate candidates
the predictor cannot resolve are its false positives. The same spacing
sets limits §5 reports separately: the 54.0\% derangement chance floor
on multi-candidate compounds (§5.5) and the ceiling neither the HOSE
lookup nor the GNN escapes (§5.4). Sharper predictors, not better
search, would move them; they exist: \textasciitilde1 ppm
message-passing ensembles\textsuperscript{36}, a DFT-coupled graph
network at 0.94 ppm\textsuperscript{37}, a community
benchmark\textsuperscript{38} and curated shift sets\textsuperscript{39}
for calibration.

The scheme is an analogy to DP4\textsuperscript{15,16} and NMR
crystallography\textsuperscript{17,18}, not an equivalence: we forgo
DFT's ≈1--2 ppm accuracy and its calibrated, probabilistic error model
for a predictor that runs in seconds.

\hypertarget{result}{%
\subsubsection{5.2 Result}\label{result}}

We ran the verifier over \textbf{every compound in the benchmark}. The
first column below is the 60-compound arm (v3 + v2-control) on which
§5.3--§5.5 build; the second is the full benchmark. All 373 candidates
were forward-predicted, so nothing here is a lower bound; the
self-ranking row re-derives Table 2 from §4's output.

\needspace{16\baselineskip}

\textbf{Table 7. Forward-verification decomposition}, original arm and
full benchmark.

\begin{longtable}[]{@{}
  >{\raggedright\arraybackslash}p{(\columnwidth - 4\tabcolsep) * \real{0.5312}}
  >{\raggedleft\arraybackslash}p{(\columnwidth - 4\tabcolsep) * \real{0.1979}}
  >{\raggedleft\arraybackslash}p{(\columnwidth - 4\tabcolsep) * \real{0.2708}}@{}}
\toprule\noalign{}
\begin{minipage}[b]{\linewidth}\raggedright
\end{minipage} & \begin{minipage}[b]{\linewidth}\raggedleft
60-compound arm
\end{minipage} & \begin{minipage}[b]{\linewidth}\raggedleft
\textbf{full benchmark (n=194)}
\end{minipage} \\
\midrule\noalign{}
\endhead
\bottomrule\noalign{}
\endlastfoot
generation recall & 19/60 (32\%) & \textbf{65/194 (34\%)} \\
top-1, solver self-ranking & 14/60 (23\%) & 55/194 (28\%) \\
top-1, \textbf{forward-verified re-ranking} & 16/60 (27\%) &
\textbf{58/194 (30\%)} \\
& & \\
\emph{conditional on recall} --- self-ranking & 14/19 (74\%) & 55/65
(\textbf{85\%}) \\
\emph{conditional on recall} --- \textbf{forward-verification} & 16/19
(\textbf{84\%}) & 58/65 (\textbf{89\%}) \\
\ldots single-candidate compounds within that set & 6/19 & 28/65 \\
& & \\
\emph{multi-candidate only} --- self-ranking & 8/13 (62\%) & 27/37
(73\%) \\
\emph{multi-candidate only} --- \textbf{forward-verification} & 10/13
(77\%) & \textbf{30/37 (81\%)} \\
\end{longtable}

\textbf{When the true structure is among the candidates, forward
verification selects it in 58 of 65 cases (89\%)}. Of the 65, \textbf{28
had a single candidate}, which any ranker scores by construction, so
that pooled figure is not the verifier's margin over chance. On the
\textbf{37 where a choice existed}, forward-verification gets
\textbf{30/37 (81\%)} against a 54.0\% derangement floor (§5.5) ---
\textbf{+27.1 points}; self-ranking \textbf{27/37 (73\%)}; we treat the
37 as the denominator that measures verification.

The margin over self-ranking remains small and unresolved: seven
compounds gained, four lost, McNemar exact \textbf{p=0.55}. We do not
claim it. The load-bearing claim is that precision is high \textbf{in
absolute terms} while recall binds: top-1 moves only 28\%→30\% because
the true structure was never proposed for 129 of 194 compounds, which no
re-ranking can repair. The verifier converts recall almost completely on
\emph{simple} targets (50/53, \textbf{94\%}) and is \textbf{exactly
tied} with self-ranking on \emph{complex} ones (8/12, 67\%). Elucidation
factorises into two near-independent levers: the \textbf{verifier is
already strong (89\%)}; the \textbf{generator (34\% recall) is the wall}
(Fig. 1).

\hypertarget{generate-wide-testing-the-recipe}{%
\subsubsection{5.3 Generate-wide: testing the
recipe}\label{generate-wide-testing-the-recipe}}

The decomposition implies a recipe: \emph{generate wide, verify by
forward prediction} --- a chemistry analog of self-consistency
sampling.\textsuperscript{40} Ten solver agents proposed up to
\textbf{six regiochemistry-aware candidates per compound}, pooled with
the originals and re-ranked as before.

\needspace{9\baselineskip}

\textbf{Table 8. Generate-wide vs original}, on the 60-compound arm
only, since wide generation was run there; the ``original'' column is
Table 7's first.

\begin{longtable}[]{@{}
  >{\raggedright\arraybackslash}p{(\columnwidth - 4\tabcolsep) * \real{0.5258}}
  >{\raggedleft\arraybackslash}p{(\columnwidth - 4\tabcolsep) * \real{0.3093}}
  >{\raggedleft\arraybackslash}p{(\columnwidth - 4\tabcolsep) * \real{0.1649}}@{}}
\toprule\noalign{}
\begin{minipage}[b]{\linewidth}\raggedright
\end{minipage} & \begin{minipage}[b]{\linewidth}\raggedleft
original (60-compound arm)
\end{minipage} & \begin{minipage}[b]{\linewidth}\raggedleft
generate-wide
\end{minipage} \\
\midrule\noalign{}
\endhead
\bottomrule\noalign{}
\endlastfoot
recall (true structure among candidates) & 32\% & \textbf{42\%} \\
forward-verified top-1 & 27\% & \textbf{30\%} \\
verification precision (conditional on recall) & 84\% & 72\% \\
\end{longtable}

Wide generation lifts recall +10 points (32\%→42\%, 19/60→25/60) and
top-1 +7 over the self-ranking baseline (23\%→30\%, 14/60→18/60),
i.e.~+3 over the forward-verified top-1 (27\%→30\%, 16/60→18/60, Table
8), on the same 60 compounds (Fig. 7) with no training.

\begin{figure}
\centering
\includegraphics[width=3.3in,height=\textheight]{docs/figures/fig3_method.png}
\caption{The forward-verification inference ladder on the 60-compound
arm: solver self-ranking → + forward-verification → + generate-wide, at
23\% / 27\% / 30\% top-1. Each rung is a training-free change to
inference alone, directional rather than individually resolved at this
sample size (§5.3).}
\end{figure}

\textbf{These top-1 differences are directional, not statistically
resolved at n=60.} The stages are \textbf{not} nested: self-ranking to
generate-wide gains seven compounds and loses three, McNemar exact
\textbf{p=0.34} (b=3, c=7). The recall gain is the better-supported
effect: the ladder shows the \emph{mechanism} works and recall is the
movable factor, not how large the top-1 gain is.

The ladder does not extend. \textbf{(i) Recall plateaus at 42\%}, on
large polycyclic targets rather than exotic ones --- \textbf{≥4 rings in
3.1\% of recalled compounds against 38.0\% of missed ones}. \textbf{(ii)
Verification precision falls 84\%→72\%} as more near-degenerate
regioisomers enter the pool. This recall/precision tension is the honest
ceiling of the training-free approach: closing the gap requires sharper
verification or 2D-NMR constraints, not merely more candidates.

\hypertarget{non-llm-verifiers-a-deterministic-lookup-and-a-learned-model}{%
\subsubsection{5.4 Non-LLM verifiers: a deterministic lookup and a
learned
model}\label{non-llm-verifiers-a-deterministic-lookup-and-a-learned-model}}

§5.3 suggests an obvious fix: a non-LLM ¹³C predictor in the verifier
slot. We tested a HOSE-code\textsuperscript{41}-style lookup and a small
message-passing GNN, both trained on the same nmrshiftdb2
dump\textsuperscript{42} and applied to the same §5.2 candidate sets, so
that only the predictor changes. The GNN is deliberately modest; sharper
purpose-built ¹³C models
exist.{[}\textsuperscript{36}{]}\textsuperscript{37}\textsuperscript{38}

\needspace{10\baselineskip}

\textbf{Table 9. Verifier comparison, conditional on recall.}

\begin{longtable}[]{@{}
  >{\raggedright\arraybackslash}p{(\columnwidth - 6\tabcolsep) * \real{0.3125}}
  >{\raggedleft\arraybackslash}p{(\columnwidth - 6\tabcolsep) * \real{0.2604}}
  >{\raggedleft\arraybackslash}p{(\columnwidth - 6\tabcolsep) * \real{0.2396}}
  >{\raggedleft\arraybackslash}p{(\columnwidth - 6\tabcolsep) * \real{0.1875}}@{}}
\toprule\noalign{}
\begin{minipage}[b]{\linewidth}\raggedright
verifier
\end{minipage} & \begin{minipage}[b]{\linewidth}\raggedleft
60-compound arm (n=19)
\end{minipage} & \begin{minipage}[b]{\linewidth}\raggedleft
full benchmark (n=65)
\end{minipage} & \begin{minipage}[b]{\linewidth}\raggedleft
held-out ¹³C MAE
\end{minipage} \\
\midrule\noalign{}
\endhead
\bottomrule\noalign{}
\endlastfoot
solver self-ranking & 14/19 (74\%) & 55/65 (85\%) & --- \\
deterministic HOSE \emph{lookup} & 14/19 (74\%) & 55/65 (85\%) & 3.23
ppm \\
\textbf{learned GNN (same data)} & 16/19 (84\%) & \textbf{59/65 (91\%)}
& \textbf{1.70 ppm} \\
LLM forward-verifier (§5.2) & 16/19 (84\%) & 58/65 (89\%) & --- \\
\end{longtable}

The lookup ties the solver's own score at both sample sizes (Table 9).
Coverage explains it: of the 6,360 candidate carbons, only \textbf{2\%}
match a training environment at r=4 and \textbf{71\%} resolve at r≤2 or
the hybridisation prior, too coarse to separate regioisomers.

The GNN matches and slightly exceeds the LLM verifier, 59/65 against
58/65 (Fig. S6), though not significantly: \emph{suggestive and
directional}, consistent with the lookup's failure being substantially
\textbf{method} rather than coverage alone, but not established. Neither
predictor reaches the near-degenerate-regioisomer precision ceiling
(§5.3/§5.5), where DFT-level accuracy or orthogonal 2D-NMR constraints
remain the genuine fix.

Leakage is slight --- no benchmark answer appears in nmrshiftdb2 at all
(0/373) --- and a Y-randomisation control (1,000 derangements) places
the real result above the 97.5th percentile of chance (n=19: real 84\%
vs mean 58.6\%, one-sided p\textless0.05). Like §5.6, the learned
verifier is a \textbf{trained complement}, reported outside the
training-free protocol.

\hypertarget{negative-control}{%
\subsubsection{5.5 Negative control}\label{negative-control}}

\textbf{Permutation negative control (Y-randomisation analog).} A
re-ranker exploiting genuine predicted-vs-observed ¹³C agreement must
lose it when the pairing is broken. We re-paired which observed ¹³C
spectrum each candidate set is scored against --- a \emph{derangement},
so no compound keeps its own spectrum --- and re-ran the verifier 1,000
times. Conditional-on-recall precision falls from the true
\textbf{89.2\% (58/65)} to a permuted mean of \textbf{73.8\%} (95\%
range 66.2--81.5\%; one-sided empirical p=0.001, two-sided p=0.002). The
verifier acts on real spectral agreement, not a candidate-list artefact.
That floor is high because 28 of the 65 recall-positive compounds carry
a single scorable candidate, and a compound with nothing to rank is
scored correct under \emph{every} pairing. Those compounds enter the
permuted score as guaranteed hits while carrying no verification signal,
so the control was partly measuring the composition of the
recall-positive set. On the \textbf{37} multi-candidate compounds ---
the set §5.2 calls the one that measures verification --- the floor
falls to \textbf{54.0\%} and the real value is \textbf{81.1\% (30/37)},
a margin of \textbf{+27.1 points} rather than +15.4, at the same
one-sided p=0.001. We report the restricted figure as the verifier's
margin over chance and the pooled one for continuity with §5.2's
denominator.

\textbf{Confidence calibration (a negative result).} Ranking the 138
multi-candidate compounds by chamfer margin and answering only the
most-confident fraction leaves top-1 flat and non-monotonic with
coverage (22\% at full coverage, 24\% at 75\%, 28\% at 50\%, 24\% at
25\%). Single-candidate compounds, which have no margin, must be
excluded; retaining them in an earlier analysis produced a spurious
``improvement''. We therefore do \textbf{not} claim the verifier
distance as a calibrated abstention gauge (§5.4).

\hypertarget{is-the-recall-wall-task-intrinsic-a-trained-generator-probe}{%
\subsubsection{5.6 Is the recall wall task-intrinsic? A
trained-generator
probe}\label{is-the-recall-wall-task-intrinsic-a-trained-generator-probe}}

The ceiling above is a property of \emph{training-free LLM elicitation};
whether the \textasciitilde42\% recall plateau is intrinsic to 1D-data
elucidation is a separate question. We probed it by substituting a small
(\textasciitilde16M-parameter) ¹H/¹³C→SMILES transformer for the
enumerator; its candidates were pooled with Claude's and re-ranked by
the same verifiers.

On the 194-compound benchmark the true structure enters the candidate
pool for \textbf{54.1\%} of compounds, versus 41.8\% for scaffold
enumeration and 33.5\% for Claude alone. Where enumeration's
near-degenerate regioisomers \emph{collapse} the HOSE verifier (top-1
28.4\%→16.0\%, the §5.3 precision-loss mechanism), the generator's
formula-correct, ¹³C-separable candidates \textbf{convert}: top-1 rises
\textbf{28.4\%→35.1\%} (McNemar exact p=0.015; +6.7 points, 95\% CI
{[}+2.1, +11.9{]}; Fig. S5). The family this is corrected over is the
probe's own three arms --- Claude alone, plus enumeration, plus
generator --- where it survives Holm (0.029) and Benjamini--Hochberg
(0.022). It would not survive a correction taken over every p-value in
the paper, and we name the family rather than leave the scope to be
guessed. With the LLM forward-verifier on the 60 forward-verify
compounds, recall rises 42\%→\textbf{57\%} (34/60) and top-1 reaches
\textbf{47\%} (28/60), at 82\% precision conditional on recall (28/34).
Those two figures are a \textbf{re-run, not a reproduction}: the earlier
pass reported 41\% top-1 at 73\% precision but deposited no outputs, so
we do not stand behind it, and the re-run lands above it.

Two controls guard against memorisation: the fine-tuning split removes
every benchmark InChIKey-14 (train∩benchmark = 0, val∩benchmark = 0),
with \textbf{none of the 40 newly-recovered compounds} in either
training stage (0/40 in simulated pretraining, 0/40 in fine-tuning); and
the simulated-pretrained model alone recovers \textbf{0/248} benchmark
structures zero-shot, rising to 25\% only after IRexp fine-tuning. The
IRexp data, not the architecture, is the active ingredient, so the
recall ceiling is \textbf{elicitation-specific, not task-intrinsic}. We
report this as a probe, not part of the headline training-free protocol.

Purpose-trained systems make the point more strongly: NMR-Solver reaches
52.9\% top-1 on experimental literature ¹H/¹³C with the formula
supplied\textsuperscript{21}, and a purpose-trained IR transformer
63.8\% top-1 on experimental NIST spectra\textsuperscript{14}. Whatever
bounds a training-free LLM at 28--30\%, it is plainly not the task. The
public split \texttt{irexp\_\hspace{0pt}release/\hspace{0pt}train} does
\emph{not} hold out the benchmark (117/200 InChIKey-14 overlap), so
downstream users must de-leak; see Data and code availability.

\begin{center}\rule{0.5\linewidth}{0.5pt}\end{center}

\hypertarget{discussion}{%
\subsection{6. Discussion}\label{discussion}}

The Claude-family models we tested are reliable \textbf{scaffold-level}
elucidators and good \textbf{verifiers} (89\% conditional on recall);
exact top-1 is throttled by candidate recall and by the regiochemical
underdetermination intrinsic to 1D NMR. The contribution is a
\textbf{diagnosis with a bounded, training-free improvement attached},
not a method that solves the task.

The diagnosis held under every perturbation we were able to apply ---
four Claude models, a second chemical domain, four verifiers --- and
reproduces inside a single application domain (§4.5): the wall is
\textbf{generation recall}, not verification, consistent with the
bottleneck being structural rather than incidental. Whether it holds
outside the Claude family is the open question, not a settled one (§7).
We measured the improvement's ceiling rather than projecting past it:
forward verification moves top-1 from 28\% to 30\% across the whole
benchmark (§5.2), while recall plateaus at 42\% (§5.3).

This reframes the engineering problem. The durable levers are not a
bespoke spectra→structure model --- a target the frontier already meets
at the scaffold level, and one that ages out each model generation ---
but open, hard, honestly scored benchmarks and inference-time
scaffolding that needs no training. Our trained-generator probe (§5.6)
shows the recall wall is \emph{elicitation-specific, not
task-intrinsic}, and the gain comes entirely from the released data
(0→25\% with fine-tuning, 0/248 without): the \textbf{open dataset}, not
a bespoke architecture, does the work when training helps.

\textbf{Protocol is a lever of the same order as capability}: a number
quoted without its protocol is uninterpretable. We do \emph{not} claim
protocol dominates capability --- the model axis spans 0\% to 46\% top-1
on the fixed 24-compound subset (Table 3), wider than the 5\%→15\%
protocol effect in §4.3, and the two are measured on different sets. In
molecular property prediction, a systematic benchmark reports that
learned, pretrained molecular embeddings rarely outperform classical
ECFP fingerprints once evaluation is controlled\textsuperscript{43}; we
read our forward-verification recipe (§5) in that light, and draw the
parallel narrowly.

Two findings transfer to practice: solve each problem in bounded,
frequently-reset contexts with tool access (5\%→15\%, 1/20→3/20,
directional at n=20, McNemar p≥0.25, §4.3); and use
forward-predicted-vs-observed ¹³C agreement to \emph{re-rank}
candidates, not to decide whether to trust the winner. Match distance
failed as an absolute confidence gauge in our selective-prediction test
(§5.5, a reported null result), so it does not support abstention
thresholds.

\begin{center}\rule{0.5\linewidth}{0.5pt}\end{center}

\hypertarget{limitations}{%
\subsection{7. Limitations}\label{limitations}}

Several limitations of earlier drafts are now resolved; we state what
remains plainly.

\emph{Resolved.} \textbf{Extraction noise:} an RDKit self-consistency
audit finds \textbf{57/60 ground truths spectrally clean}, with metrics
unchanged on that subset. \textbf{Verifier sample size:} forward
verification now runs over \textbf{all 194 compounds}, so the
recall-conditional claim rests on n=65 rather than n=19 (§5.2).
\textbf{Verifier precision:} precision degrades as near-degenerate
regioisomers accumulate (§5.3), so forward-match distance is a strong
\emph{re-ranker} but a soft \emph{confidence} gauge. \textbf{Non-LLM
verifiers} (§5.4): a HOSE-code \emph{lookup} does \textbf{not} beat the
LLM verifier (85\% vs 89\% conditional at n=65) and a \emph{learned} GNN
reaches 91\% --- directional but not statistically resolved at n=65, so
it is suggestive that the deterministic failure was method rather than
coverage alone, not a demonstration of it. Beyond all of them lies the
near-degenerate-regiochemistry precision ceiling, where DFT-level
accuracy or 2D-NMR constraints are the genuine fix. \textbf{Projection:}
§5.3 now measures what an earlier draft estimated: top-1 30\%, recall
42\%, below the estimate it replaces, reported as found rather than
adjusted.

\emph{Independence checks.} Scoring is mechanical RDKit, not LLM-judged,
and all solver/verifier runs were transcript-audited \emph{at generation
time} for zero web and zero ground-truth access. That audit is the one
part of the pipeline a reader \textbf{cannot re-verify from the
release}: the committed artefacts are the parsed per-compound
predictions, and the transcripts are not deposited (they are available
on request). The \emph{outcome} of closed-book solving is separately
testable: the formula-only control (§4.6) removes the spectra and
accuracy collapses, and a permutation negative control (§5.5) collapses
conditional-on-recall precision from 89.2\% to a chance mean of 73.8\%
(two-sided p=0.002 on the full benchmark), as a correctly isolated blind
evaluation requires. A second Claude model (Sonnet) was comparably
recall-bound on a 12-compound subset (recall 33\% vs Opus 42\% on those
compounds), consistent with the recall bound being a property of the
task rather than of one model --- but within the Claude family, so it
speaks to model-instance robustness, not cross-vendor generality.

\emph{Remaining.} \textbf{(i) Pretraining contamination is bounded, not
excluded.} Every benchmark compound was mined from open-access
literature, so a frontier model may have encountered it in training. The
formula-only control (§4.6) bounds how much of the headline number pure
recall can explain: masking the spectra drops top-1 from 23\% to 5\%,
perfectly nested, McNemar exact p=0.001. It examines the three
formula-only successes individually rather than collectively --- only
one (an unlisted silyl aryl sulfonate) has a genuinely determining
composition, while \emph{N}-tosyl-leucine and {[}6{]}-paradol are
catalogued compounds for which recall is a live explanation --- which
leaves the bound unchanged. A second control finds accuracy flat in
source publication year across all 194 compounds (r=−0.007), the
size-adjusted older-versus-newer difference bounded at −5.1 points, 95\%
CI {[}−17.2, +7.0{]}. Neither control is randomised: publication year is
observational, and the formula-only arm cannot separate memorisation
from inference on a near-determining formula. Nor is the recency test
anchored to a \emph{known} training cutoff, since the harness does not
disclose one (§8); we test recency as a continuous proxy. A replication
restricted to compounds published after a disclosed cutoff remains open.

The sharper form of this objection is one neither control addresses, and
we state it plainly because a reader is entitled to weigh it. The prompt
carries the spectral strings \emph{exactly as printed in the source
article}, down to its own typography --- the sampler keeps a compound
only when the raw ¹H payload can be recovered verbatim from the
open-access full text (§8), so \textbf{every benchmark compound has its
spectrum present verbatim in a document the model may have read}, and
that string is a high-entropy fingerprint of the document. Masking the
spectra therefore removes the chemistry and a retrieval key together: a
collapse from 23\% to 5\% is what reading predicts, and also what
retrieval predicts. Recency cannot separate them either, for the reason
just given. The experiment that would is a spectrum meaning the same
thing and reading differently --- shifts perturbed within reporting
precision and typography normalised, so the chemistry is untouched and
no verbatim string survives.
\texttt{scripts/\hspace{0pt}jitter\_\hspace{0pt}control.py} builds that
set from the released benchmark; running it is the obvious next control
and we have not run it. It would bound verbatim retrieval only, not
recognition from approximate shift patterns.

\textbf{(ii) Human audit --- prepared but not yet run.} Two things are
unvalidated by hand: the elucidation and forward-prediction outputs, and
the \emph{recall} side of dataset extraction (whether the parser found
every IR string in every source paper). Transcription fidelity of the
records we hold is measured and high (§2.3); record-level recall is not,
and only reading papers can settle it. Solver and verifier are both
LLMs, so the one validation we cannot perform ourselves is
expert-chemist review. We have therefore \textbf{built and frozen} a
blinded, pre-registered audit package at
\texttt{data/\hspace{0pt}audit/\hspace{0pt}}, protocol at
\texttt{docs/\hspace{0pt}EXPERT\_\hspace{0pt}AUDIT\_\hspace{0pt}PROTOCOL.md}:
a difficulty-stratified 30-compound elucidation panel plus seven
ranking-only compounds, with a withheld answer key, testing what a miss
actually is (§4) and whether forward verification is a trustworthy
re-ranker (§5). \textbf{We have not yet run the panel}; until expert
results are in, those claims should be read as machine-validated (RDKit
InChIKey) but not yet human-validated. §4 has since narrowed its first
question by reporting all 137 analysable misses mechanically (76.6\%
constitutional isomers, 22.6\% scaffold-preserving positional errors),
correcting the impressionistic claim the audit was built to check; what
remains for a chemist is whether a formula-correct, scaffold-wrong
candidate is a \emph{chemically reasonable} reading of the spectra.

\textbf{(iii) Single vendor and an underpowered cross-model comparison.}
The headline (28.4\% top-1, n=194) is a single frontier model (Claude
Opus) on the full benchmark; the cross-model evidence (§4.4) is four
\textbf{Claude-family} models on a shared \textbf{24-compound} subset,
and is underpowered: the ranking is robust (strictly nested, weakest
floored at 0\%) but no adjacent gap is established at n=24, so
quantitative claims should be read as holding \textbf{for the Claude
family}. The cross-\textbf{vendor} question is no longer open: §4.7
finds verification precision exceeding generation recall for Grok 4.6,
Gemini 3.7 Flash and GPT-5.6 Sol on the 60-compound arm, as for Claude.
What remains is weaker than the original gap but real: only Claude's
inequality is interval-disjoint at n=60, the other three are
directional. Those models were reached through a coding-assistant
harness exposing reasoning-effort tiers, a decoding control our own arm
never had, so they measure a model as served rather than a bare
endpoint. The clean-clone contamination control was run for Grok alone,
and the cross-vendor arms cover 60 compounds.

\textbf{(iv) Domain-subset scope:} the battery-electrolyte case study
(§4.5) uses \emph{literature compounds bearing electrolyte-relevant
functional chemistry}, not operando or in-cell decomposition spectra, so
it demonstrates functional-class transfer of the elucidation bottleneck,
not direct assignment of authentic interphase/degradation products,
which would require a dedicated operando set.

\textbf{(v) Constitution-only scoring:} correctness is judged on
InChIKey connectivity, so a correct-constitution / wrong-stereochemistry
prediction counts as correct. Scoring the full InChIKey gives 21.1\%
top-1 (§3), 7.3 points lower, so the headline should be read as an upper
bound on full-stereochemistry accuracy. We keep constitution because 1D
NMR/IR rarely determines absolute configuration --- the 10.3\% (20/194)
of targets with a defined stereocentre would be penalised for
information the prompt never carried --- but a stereochemistry-sensitive
benchmark would need 2D/chiroptical data.

\textbf{(vi) Single-sample scoring:} each headline compound is scored
from one solver prediction set, so reported top-1/recall carry no
run-to-run (LLM-sampling) variance estimate; the bootstrap CIs reflect
compound sampling only. §5.3 pools ten independent generation passes
(recall 32\%→42\%), indicating generator stochasticity that single-pass
scoring understates.

\textbf{(vii) Chemical-space coverage:} the benchmark is drawn from
open-access organic methodology and total-synthesis literature; it is
\textbf{not} representative of organometallic/coordination compounds,
large biomolecules (peptides, oligonucleotides, oligosaccharides), or
stereochemistry-heavy targets, and our results should not be
extrapolated to them.

\begin{center}\rule{0.5\linewidth}{0.5pt}\end{center}

\hypertarget{methods}{%
\subsection{8. Methods}\label{methods}}

\textbf{Mining and resolution.} PMC-OA full text was fetched from
\texttt{s3:/\hspace{0pt}/\hspace{0pt}pmc-oa-opendata}, parsed
deterministically, and resolved with OPSIN, RDKit and SELFIES, with an
optional cached PubChem fallback.

\textbf{Benchmark and agents.} Problems were sampled from
\texttt{irexp\_\hspace{0pt}resolved}, stratified by RDKit ring analysis
to half the round per stratum, and de-duplicated across rounds by
InChIKey. A compound was admitted only if its raw ¹H payload ---
multiplicities and \emph{J} values as printed --- could be re-extracted
verbatim from the PMC-OA full text of its source article
(\texttt{benchmark\_\hspace{0pt}v2.raw\_\hspace{0pt}1h\_\hspace{0pt}for}),
which is what puts genuine coupling information in the prompt and,
unavoidably, puts the source article's exact string there too (§7). The
battery-electrolyte subset (§4.5) was drawn from the same corpus by
SMARTS filters for six electrolyte functional classes (carbonate,
sulfonyl/sulfonate, nitrile, sp³-C--F, phosphoryl, glyme/oligoether),
balanced to eight compounds per class (48 curated; 46 scored after two
yielded no parseable candidate), excluding every compound used
elsewhere, and J-enriched and spectrally validated as in the main
rounds. Solver and forward-prediction agents were independent
Claude-Opus sub-agents under the same subscription, instructed
closed-book and audited by automated transcript search for zero
web/answer access. A prediction was scored correct on RDKit
InChIKey-connectivity match; similarity used Morgan(2, 2048) Tanimoto,
and forward-verification a symmetric chamfer distance over ¹³C peak
sets. The core protocol trains no model and uses no paid API; the §5.6
generator and the §5.4 learned ¹³C verifier are the only two trained
components, reported as complements and fenced from the headline
results.

\emph{Models and versions.} All experiments used Anthropic Claude models
via the consumer subscription (claude.ai). The §4.4 comparison spanned,
in capability order, Claude Haiku, Claude Sonnet, Claude Opus and Claude
Fable 5; the headline benchmark (§4.1) and forward-verification (§5)
used Claude Opus. \texttt{docs/\hspace{0pt}MODELS.md} records, per
experiment, the model, data directory, collection date, harness and tool
access, and scoring code path, across three dated windows:

\begin{itemize}
\tightlist
\item
  \textbf{2026-06-09 to 2026-06-11} (UTC) --- every candidate structure
  behind the \textbf{headline} results: §4.1, §4.3--§4.5, and the
  candidate pools all of §5 re-ranks.
\item
  \textbf{2026-07-28} --- the formula-only contamination control (§4.6),
  re-solving the same 60 compounds with spectra masked and generating
  its own candidates; it touches Table 5 alone.
\item
  \textbf{2026-08-07} --- three collections forward-predicting ¹³C for
  candidates the June run had already produced (the §5.2 extension to
  all 194 compounds, the §5.3 coverage-gap closure, and the §5.6
  re-run); none introduced a new candidate structure, and no recall
  number in the paper changes.
\end{itemize}

One \emph{verified} number does: the §5.6 re-run supersedes a top-1
whose original forward-prediction outputs were never deposited and are
lost, disagreeing with the figure it replaces (47\% against 41\%); §5.6
reports it as a re-run rather than a reproduction.

\textbf{No model snapshot can be reported.} The consumer harness exposes
no checkpoint identifier, records nothing about which build served a
given request, and exposes no decoding parameters, which were neither
set nor recorded. \texttt{Claude\ Opus\ 4.8} is evidenced only for the
2026-06-09 pilot, so \textbf{a mid-window build change cannot be
excluded} and we do not claim the same build served the main round two
days later. Re-running reproduces the protocol distributionally rather
than exactly, and exact agreement is not offered as a target. Fixed are
the dated collection windows, the frozen per-compound outputs, and
mechanical scorers that regenerate every number from them: exact
reproducibility of scoring and analysis, not of inference
(\texttt{docs/\hspace{0pt}MODELS.md} §6).

\textbf{Reproducibility.} Every round is frozen: questions, ground-truth
answers, per-agent raw outputs, predictions and scorer outputs are
released; the sampler, scorer and forward-verification harness are
scripted end-to-end.

\begin{center}\rule{0.5\linewidth}{0.5pt}\end{center}

\hypertarget{supporting-information-figures}{%
\subsection{Supporting Information
figures}\label{supporting-information-figures}}

These are supplied as a separate Electronic Supplementary Information
document (\texttt{docs/\hspace{0pt}paper\_\hspace{0pt}esi.pdf}, built by
\texttt{scripts/\hspace{0pt}build\_\hspace{0pt}pdf.py}), as RSC
requires; they are listed here for reference.

\begin{itemize}
\tightlist
\item
  \textbf{Fig. S1}
  (\texttt{docs/\hspace{0pt}figures/\hspace{0pt}fig0\_\hspace{0pt}overview.png})
  --- study design: open multimodal data (IRexp) → blind,
  complexity-stratified benchmark → decoupled blind solving →
  forward-verification re-ranking; training-free core pipeline.
\item
  \textbf{Fig. S2}
  (\texttt{docs/\hspace{0pt}figures/\hspace{0pt}fig4\_\hspace{0pt}dataset.png})
  --- IRexp composition: IR records → NMR-paired → structure-linked →
  full IR+¹H+¹³C+structure quadruples.
\item
  \textbf{Fig. S3}
  (\texttt{docs/\hspace{0pt}figures/\hspace{0pt}fig2\_\hspace{0pt}size.png})
  --- accuracy vs molecular size (heavy-atom bucket); the monotonic
  60\%→7\% top-1 gradient.
\item
  \textbf{Fig. S4}
  (\texttt{docs/\hspace{0pt}figures/\hspace{0pt}fig6\_\hspace{0pt}electrolyte.png})
  --- top-1 and recovered accuracy on IRSpectra-Bench-Electrolyte by
  battery-electrolyte chemical class (n=46): sp³-C--F easiest (50\%),
  sulfonyl and nitrile hardest (12\%); overall 26\%/28\%.
\item
  \textbf{Fig. S5}
  (\texttt{docs/\hspace{0pt}figures/\hspace{0pt}fig\_\hspace{0pt}generator\_\hspace{0pt}probe.png})
  --- trained-generator probe (§5.6; a complement, not part of the
  training-free protocol): candidate recall and deterministic-HOSE top-1
  on the 194-compound benchmark for Claude / + scaffold enumeration / +
  trained generator. Enumeration's near-degenerate isomers collapse the
  verifier (28.4→16.0\%) while the generator's formula-correct
  candidates convert (28.4→35.1\%).
\item
  \textbf{Fig. S6}
  (\texttt{docs/\hspace{0pt}figures/\hspace{0pt}fig\_\hspace{0pt}verifier.png})
  --- learned-verifier probe (§5.4; a complement, not part of the
  training-free protocol). (A) Conditional-on-recall top-1 (n=65, whole
  benchmark) for the four verifiers: a GNN trained on the same
  nmrshiftdb2 data as the HOSE lookup reaches the LLM verifier's level
  (91\% vs 89\%) where the lookup (85\%) does not move off the solver's
  own ranking. (B) Why: held-out ¹³C MAE --- the learned model is
  \textasciitilde2× sharper (1.70 vs 3.23 ppm).
\end{itemize}

\begin{center}\rule{0.5\linewidth}{0.5pt}\end{center}

\hypertarget{data-and-code-availability}{%
\subsection{Data and code
availability}\label{data-and-code-availability}}

Everything below is in the project repository. The archival deposit ---
a complete frozen snapshot of dataset, benchmark, answer keys,
predictions, scripts, figure regeneration and the expert-audit package
--- will carry DOI \textbf{{[}TODO: 10.5281/zenodo.XXXXXXX --- mint on
submission{]}}; GitHub is the development mirror and the Zenodo record
the citable version.

\begin{longtable}[]{@{}
  >{\raggedright\arraybackslash}p{(\columnwidth - 4\tabcolsep) * \real{0.2292}}
  >{\raggedright\arraybackslash}p{(\columnwidth - 4\tabcolsep) * \real{0.3333}}
  >{\raggedright\arraybackslash}p{(\columnwidth - 4\tabcolsep) * \real{0.4375}}@{}}
\toprule\noalign{}
\begin{minipage}[b]{\linewidth}\raggedright
component
\end{minipage} & \begin{minipage}[b]{\linewidth}\raggedright
data
\end{minipage} & \begin{minipage}[b]{\linewidth}\raggedright
regenerated by
\end{minipage} \\
\midrule\noalign{}
\endhead
\bottomrule\noalign{}
\endlastfoot
IRexp and its structure-complete split &
\texttt{data/\hspace{0pt}irexp/\hspace{0pt}},
\texttt{data/\hspace{0pt}irexp\_\hspace{0pt}resolved/\hspace{0pt}} &
\texttt{spectro\_\hspace{0pt}scraper/\hspace{0pt}} (mining pipeline) \\
benchmark rounds and the within-compound control &
\texttt{data/\hspace{0pt}benchmark*/\hspace{0pt}} &
\texttt{scripts/\hspace{0pt}benchmark\_\hspace{0pt}v2.py} \\
ground-truth integrity audit &
\texttt{data/\hspace{0pt}benchmark*/\hspace{0pt}clean\_\hspace{0pt}qids.json}
& \texttt{scripts/\hspace{0pt}validate\_\hspace{0pt}benchmark.py} \\
headline accuracy (Table 2) & --- &
\texttt{scripts/\hspace{0pt}score\_\hspace{0pt}main.py} \\
battery-electrolyte subset (§4.5) &
\texttt{data/\hspace{0pt}benchmark\_\hspace{0pt}electrolyte/\hspace{0pt}}
&
\texttt{scripts/\hspace{0pt}build\_\hspace{0pt}electrolyte\_\hspace{0pt}bench.py},
\texttt{scripts/\hspace{0pt}score\_\hspace{0pt}electrolyte.py} \\
formula-only contamination control (§4.6) &
\texttt{data/\hspace{0pt}modality/\hspace{0pt}} &
\texttt{scripts/\hspace{0pt}modality\_\hspace{0pt}ablation.py} \\
publication-recency control (§4.6) &
\texttt{data/\hspace{0pt}audit/\hspace{0pt}recency\_\hspace{0pt}control.json}
& \texttt{scripts/\hspace{0pt}contamination\_\hspace{0pt}recency.py} \\
forward-verification, original arm (§5.2) &
\texttt{data/\hspace{0pt}fverify/\hspace{0pt}} &
\texttt{scripts/\hspace{0pt}forward\_\hspace{0pt}verify.py} \\
\ldots extended to all 194 compounds &
\texttt{data/\hspace{0pt}fverify\_\hspace{0pt}main/\hspace{0pt}} &
\texttt{scripts/\hspace{0pt}forward\_\hspace{0pt}verify\_\hspace{0pt}main.py},
\texttt{scripts/\hspace{0pt}forward\_\hspace{0pt}verify\_\hspace{0pt}all.py} \\
generate-wide arm (§5.3) & \texttt{data/\hspace{0pt}gw/\hspace{0pt}},
\texttt{data/\hspace{0pt}fverify2/\hspace{0pt}} &
\texttt{scripts/\hspace{0pt}score\_\hspace{0pt}generate\_\hspace{0pt}wide.py},
\texttt{scripts/\hspace{0pt}ladder\_\hspace{0pt}significance.py} \\
\ldots its coverage gap, closed &
\texttt{data/\hspace{0pt}fverify\_\hspace{0pt}gw/\hspace{0pt}} &
\texttt{scripts/\hspace{0pt}forward\_\hspace{0pt}verify\_\hspace{0pt}gw.py} \\
non-LLM verifier comparison (Table 9) &
\texttt{data/\hspace{0pt}fverify/\hspace{0pt}hose\_\hspace{0pt}results.txt},
\texttt{data/\hspace{0pt}fverify/\hspace{0pt}verifier\_\hspace{0pt}table\_\hspace{0pt}results.txt}
& \texttt{scripts/\hspace{0pt}hose\_\hspace{0pt}predict.py}
(incl.~\texttt{coverage}),
\texttt{scripts/\hspace{0pt}verifier\_\hspace{0pt}table.py},
\texttt{scripts/\hspace{0pt}verifier\_\hspace{0pt}leakage.py} \\
negative control and selective prediction (§5.5) & --- &
\texttt{scripts/\hspace{0pt}verifier\_\hspace{0pt}diagnostics.py} \\
what a miss is, and isomer separability (§4.1, §5.1) & --- &
\texttt{scripts/\hspace{0pt}analyze\_\hspace{0pt}misses.py},
\texttt{scripts/\hspace{0pt}isomer\_\hspace{0pt}separability.py} \\
difficulty-threshold sensitivity (§3) & --- &
\texttt{scripts/\hspace{0pt}difficulty\_\hspace{0pt}sensitivity.py} \\
licence-pool split (§2.1) & --- &
\texttt{scripts/\hspace{0pt}split\_\hspace{0pt}license\_\hspace{0pt}pools.py} \\
recall headroom and scaffold enumeration & --- &
\texttt{scripts/\hspace{0pt}analyze\_\hspace{0pt}recall\_\hspace{0pt}headroom.py},
\texttt{scripts/\hspace{0pt}enumerate\_\hspace{0pt}isomers.py},
\texttt{scripts/\hspace{0pt}closing\_\hspace{0pt}the\_\hspace{0pt}gap.py} \\
blinded expert-audit package (§7) &
\texttt{data/\hspace{0pt}audit/\hspace{0pt}} &
\texttt{scripts/\hspace{0pt}make\_\hspace{0pt}audit\_\hspace{0pt}sample.py} \\
manuscript integrity gates & --- &
\texttt{scripts/\hspace{0pt}check\_\hspace{0pt}manuscript.py},
\texttt{scripts/\hspace{0pt}verify\_\hspace{0pt}statistics.py} \\
\end{longtable}

\textbf{Trained complements}, both fenced from the training-free
protocol. The §5.6 generator ships as a self-contained bundle at
\texttt{contrib/\hspace{0pt}generator\_\hspace{0pt}probe/\hspace{0pt}}
--- candidates, the de-leaked split with its InChIKey-14 manifest, the
verification scripts
(\texttt{scripts/\hspace{0pt}closing\_\hspace{0pt}the\_\hspace{0pt}gap\_\hspace{0pt}gen.py},
\texttt{scripts/\hspace{0pt}forward\_\hspace{0pt}verify\_\hspace{0pt}gen.py},
\texttt{scripts/\hspace{0pt}verify\_\hspace{0pt}leakage\_\hspace{0pt}exact40.py})
and the blind forward-prediction outputs behind its verified top-1
(\texttt{data/\hspace{0pt}fverify\_\hspace{0pt}gen/\hspace{0pt}}), so
its scorer runs with no missing predictions. The §5.4 learned verifier
ships \texttt{scripts/\hspace{0pt}gnn\_\hspace{0pt}predict.py} (extract
/ train / score / control), the trained model
\texttt{data/\hspace{0pt}nmrshiftdb/\hspace{0pt}gnn\_\hspace{0pt}c13.pt},
per-compound results and both leakage checks
(\texttt{data/\hspace{0pt}fverify/\hspace{0pt}gnn\_\hspace{0pt}results.txt}),
and the write-up
\texttt{docs/\hspace{0pt}VERIFIER\_\hspace{0pt}PROBE.md}. Model
checkpoints are deposited on Zenodo. Both trained predictors of §5.4
need the nmrshiftdb2 dump, which we cannot redistribute; the README
gives the one-line fetch.

Companion technical notes: \texttt{docs/\hspace{0pt}BENCHMARK.md},
\texttt{docs/\hspace{0pt}FORWARD\_\hspace{0pt}VERIFY.md},
\texttt{docs/\hspace{0pt}MODALITY\_\hspace{0pt}ABLATION.md},
\texttt{docs/\hspace{0pt}EXPERT\_\hspace{0pt}AUDIT\_\hspace{0pt}PROTOCOL.md},
\texttt{docs/\hspace{0pt}MODELS.md}.

\textbf{Two cautions for re-users.} The public
\texttt{irexp\_\hspace{0pt}release/\hspace{0pt}train} split does
\textbf{not} hold out IRSpectra-Bench --- it overlaps by 117/200
InChIKey-14 --- so de-leak with
\texttt{contrib/\hspace{0pt}generator\_\hspace{0pt}probe/\hspace{0pt}build\_\hspace{0pt}exp\_\hspace{0pt}manifest.py}
before training anything that will be evaluated on the benchmark. And
the held-out answer keys
\texttt{data/\hspace{0pt}audit/\hspace{0pt}key.jsonl} and
\texttt{data/\hspace{0pt}modality/\hspace{0pt}key.json} are deposited
but flagged \emph{withhold from blinded reviewers}.

\hypertarget{licensing-and-attribution}{%
\subsection{Licensing and attribution}\label{licensing-and-attribution}}

IRexp is derived from two open-access source pools and redistributed
under terms compatible with each (see \texttt{data/\hspace{0pt}NOTICE}).
\textbf{(a) Redistribution:} we release only \emph{extracted numeric
data} --- IR band lists, ¹H/¹³C shift lists, and resolved structures
(SMILES/SELFIES/InChIKey) --- plus each record's source DOI/accession,
not source full text, figures, or PDFs. \textbf{(b) Two separable
pools:} 119,345 records derive from the PMC Open-Access Subset
(\textbf{CC-BY-4.0}) and 1,888 from the Chemotion RADAR4Chem FT-IR
deposit (\textbf{CC-BY-SA-4.0}). The two are separable losslessly by
\texttt{source\_\hspace{0pt}doi} --- Chemotion records carry the
\texttt{10.22000} prefix --- and
\texttt{scripts/\hspace{0pt}split\_\hspace{0pt}license\_\hspace{0pt}pools.py}
materialises them as two files with each record stamped
\texttt{license}, so users may take the CC-BY pool alone; any combined
or Chemotion-derived release carries CC-BY-SA-4.0 to honour the
ShareAlike term. Code is released under the MIT License. \textbf{(c)
Attribution:} re-users must cite this dataset (Zenodo DOI above) and
attribute the original publications via each record's
\texttt{source\_\hspace{0pt}doi}.

\hypertarget{author-contributions}{%
\subsection{Author Contributions}\label{author-contributions}}

\textbf{I.Y.:} conceptualization, methodology, software, formal
analysis, investigation, data curation, visualization, writing ---
original draft. \textbf{R.S.:} methodology, software, formal analysis,
investigation, validation (trained-generator and learned-verifier
probes, §5.4 and §5.6), writing --- review and editing.
\textbf{R.A.V.-H.:} conceptualization, methodology, supervision, writing
--- review and editing.

\hypertarget{conflicts-of-interest}{%
\subsection{Conflicts of Interest}\label{conflicts-of-interest}}

The authors declare no competing interests.

\hypertarget{use-of-ai-tools}{%
\subsection{Use of AI tools}\label{use-of-ai-tools}}

This work studies a large language model, and LLMs were also used as
instruments and as writing aids; we state both roles explicitly.
\textbf{As an object of study:} all reported elucidation,
forward-prediction and verification results were produced by Claude
models invoked under the protocol of §3 and §8 --- these are the
measurements the paper reports, not assistance in producing it.
\textbf{As a writing aid:} the authors used an LLM-based coding
assistant for figure generation, analysis scripting, manuscript
copy-editing, and for the internal review pass that produced several of
the corrections recorded in the repository history. \textbf{No text,
number, figure or citation in this manuscript was accepted without
author verification against the released data and code}; every
quantitative claim regenerates from the scripts in
\texttt{scripts/\hspace{0pt}}. The authors take full responsibility for
the content.

\hypertarget{acknowledgements}{%
\subsection{Acknowledgements}\label{acknowledgements}}

\emph{(To be completed before submission: funding sources,
compute/infrastructure, and any individual acknowledgements. ---
AUTHORS)}

\hypertarget{references}{%
\subsection{References}\label{references}}

\hypertarget{refs}{}
\begin{CSLReferences}{0}{0}
\leavevmode\vadjust pre{\hypertarget{ref-chacko2024spectro}{}}%
\CSLLeftMargin{1 }%
\CSLRightInline{E. Chacko, R. Sondhi, A. Praveen, K. L. Luska and R. A.
Vargas-Hernández, \emph{ChemRxiv}, 2024, preprint, DOI:
\href{https://doi.org/10.26434/chemrxiv-2024-37v2j}{10.26434/chemrxiv-2024-37v2j}.}

\leavevmode\vadjust pre{\hypertarget{ref-kamber2026chemist}{}}%
\CSLLeftMargin{2 }%
\CSLRightInline{D. Kamber,
\emph{\href{https://anthropic.com/research/making-claude-a-chemist}{Making
Claude a chemist: How Claude performs on NMR prediction and structure
elucidation}}, Anthropic, 2026.}

\leavevmode\vadjust pre{\hypertarget{ref-hu2024multitask}{}}%
\CSLLeftMargin{3 }%
\CSLRightInline{\href{https://doi.org/10.1021/acscentsci.4c01132}{F. Hu,
M. S. Chen, G. M. Rotskoff, M. W. Kanan and T. E. Markland, \emph{ACS
Cent. Sci.}, 2024, \textbf{10}, 2162--2170}.}

\leavevmode\vadjust pre{\hypertarget{ref-yang2026nmrtrans}{}}%
\CSLLeftMargin{4 }%
\CSLRightInline{L. Yang, Z. Yang, J. Xie, Y. Wang, B. Gao, T. Fu, X.
Wei, J. Sun, J. Wu, C. He, Y. Li and Q. Gu, \emph{arXiv}, 2026,
preprint, DOI:
\href{https://doi.org/10.48550/arXiv.2602.10158}{10.48550/arXiv.2602.10158}.}

\leavevmode\vadjust pre{\hypertarget{ref-ottomano2025nmiracle}{}}%
\CSLLeftMargin{5 }%
\CSLRightInline{F. Ottomano, Y. Li and A. M. Ganose, \emph{arXiv}, 2025,
preprint, DOI:
\href{https://doi.org/10.48550/arXiv.2512.19733}{10.48550/arXiv.2512.19733}.}

\leavevmode\vadjust pre{\hypertarget{ref-mbran2024chemcrow}{}}%
\CSLLeftMargin{6 }%
\CSLRightInline{\href{https://doi.org/10.1038/s42256-024-00832-8}{A. M.
Bran, S. Cox, O. Schilter, C. Baldassari, A. D. White and P. Schwaller,
\emph{Nat. Mach. Intell.}, 2024, \textbf{6}, 525--535}.}

\leavevmode\vadjust pre{\hypertarget{ref-boiko2023coscientist}{}}%
\CSLLeftMargin{7 }%
\CSLRightInline{\href{https://doi.org/10.1038/s41586-023-06792-0}{D. A.
Boiko, R. MacKnight, B. Kline and G. Gomes, \emph{Nature}, 2023,
\textbf{624}, 570--578}.}

\leavevmode\vadjust pre{\hypertarget{ref-su2025spectrallm}{}}%
\CSLLeftMargin{8 }%
\CSLRightInline{Y. Su, J. Chen, Z. Jiang, Z. Zhong, L. Wang, Q. Liu and
Z. Zhang, \emph{arXiv}, 2025, preprint, DOI:
\href{https://doi.org/10.48550/arXiv.2508.08441}{10.48550/arXiv.2508.08441}.}

\leavevmode\vadjust pre{\hypertarget{ref-shen2025specmol}{}}%
\CSLLeftMargin{9 }%
\CSLRightInline{S. Shen, J. Xie, Z. Yang, A. Zhang, S. Sun, B. Gao, T.
Fu, B. Qi and Y. Li, \emph{arXiv}, 2025, preprint, DOI:
\href{https://doi.org/10.48550/arXiv.2509.21861}{10.48550/arXiv.2509.21861}.}

\leavevmode\vadjust pre{\hypertarget{ref-zhuang2025treesearch}{}}%
\CSLLeftMargin{10 }%
\CSLRightInline{\href{https://aclanthology.org/2025.acl-long.1100/}{X.
Zhuang, B. Wu, J. Cui, K. Feng, X. Li, H. Xing, K. Ding, Q. Zhang and H.
Chen, \emph{Proc. 63rd Annu. Meet. Assoc. Comput. Linguist. (Vol. 1:
Long Papers)}, 2025, 22561--22576}.}

\leavevmode\vadjust pre{\hypertarget{ref-guo2024molpuzzle}{}}%
\CSLLeftMargin{11 }%
\CSLRightInline{\href{https://doi.org/10.52202/079017-4281}{K. Guo, B.
Nan, Y. Zhou, T. Guo, Z. Guo, M. Surve, Z. Liang, N. V. Chawla, O. Wiest
and X. Zhang, in \emph{Advances in neural information processing systems
37 (NeurIPS 2024), datasets and benchmarks track}, 2024}.}

\leavevmode\vadjust pre{\hypertarget{ref-noh2025iragent}{}}%
\CSLLeftMargin{12 }%
\CSLRightInline{H. Noh, N. Lee, G. S. Na, K. Kim and C. Park,
\emph{arXiv}, 2025, preprint, DOI:
\href{https://doi.org/10.48550/arXiv.2508.16112}{10.48550/arXiv.2508.16112}.}

\leavevmode\vadjust pre{\hypertarget{ref-alberts2024ir}{}}%
\CSLLeftMargin{13 }%
\CSLRightInline{\href{https://doi.org/10.1038/s42004-024-01341-w}{M.
Alberts, T. Laino and A. C. Vaucher, \emph{Commun. Chem.}, 2024,
\textbf{7}, 268}.}

\leavevmode\vadjust pre{\hypertarget{ref-alberts2025benchmarks}{}}%
\CSLLeftMargin{14 }%
\CSLRightInline{\href{https://doi.org/10.1039/d5dd00131e}{M. Alberts, F.
Zipoli and T. Laino, \emph{Digital Discovery}, 2025, \textbf{4},
1936--1943}.}

\leavevmode\vadjust pre{\hypertarget{ref-smith2010dp4}{}}%
\CSLLeftMargin{15 }%
\CSLRightInline{\href{https://doi.org/10.1021/ja105035r}{S. G. Smith and
J. M. Goodman, \emph{J. Am. Chem. Soc.}, 2010, \textbf{132},
12946--12959}.}

\leavevmode\vadjust pre{\hypertarget{ref-grimblat2015dp4plus}{}}%
\CSLLeftMargin{16 }%
\CSLRightInline{\href{https://doi.org/10.1021/acs.joc.5b02396}{N.
Grimblat, M. M. Zanardi and A. M. Sarotti, \emph{J. Org. Chem.}, 2015,
\textbf{80}, 12526--12534}.}

\leavevmode\vadjust pre{\hypertarget{ref-pickard2001gipaw}{}}%
\CSLLeftMargin{17 }%
\CSLRightInline{\href{https://doi.org/10.1103/PhysRevB.63.245101}{C. J.
Pickard and F. Mauri, \emph{Phys. Rev. B}, 2001, \textbf{63}, 245101}.}

\leavevmode\vadjust pre{\hypertarget{ref-ashbrook2016nmrcryst}{}}%
\CSLLeftMargin{18 }%
\CSLRightInline{\href{https://doi.org/10.1039/c6cc02542k}{S. E. Ashbrook
and D. McKay, \emph{Chem. Commun.}, 2016, \textbf{52}, 7186--7204}.}

\leavevmode\vadjust pre{\hypertarget{ref-elyashberg2012case}{}}%
\CSLLeftMargin{19 }%
\CSLRightInline{\href{https://doi.org/10.1002/mrc.2849}{M. Elyashberg,
K. Blinov, S. Molodtsov and A. Williams, \emph{Magn. Reson. Chem.},
2012, \textbf{50}, 22--27}.}

\leavevmode\vadjust pre{\hypertarget{ref-elyashberg2015case}{}}%
\CSLLeftMargin{20 }%
\CSLRightInline{M. E. Elyashberg and A. J. Williams,
\emph{\href{https://doi.org/10.1007/978-3-662-46402-1}{Computer-based
structure elucidation from spectral data}}, Springer Berlin Heidelberg,
2015.}

\leavevmode\vadjust pre{\hypertarget{ref-jin2025nmrsolver}{}}%
\CSLLeftMargin{21 }%
\CSLRightInline{Y. Jin, J.-J. Wang, F. Xu, X. Ji, Z. Gao, L. Zhang, G.
Ke, R. Zhu and W. E, \emph{Nat. Commun.},
DOI:\href{https://doi.org/10.1038/s41467-026-71315-0}{10.1038/s41467-026-71315-0}.}

\leavevmode\vadjust pre{\hypertarget{ref-fang2026nmragent}{}}%
\CSLLeftMargin{22 }%
\CSLRightInline{Z. Fang, C. Yang, Y. Tan, Y. Zhao, F. Xu, H. Xiang, H.
Sun, H. Gao, X. Wang, W. Du, Y. Li and J. Xia, \emph{arXiv}, 2026,
preprint, DOI:
\href{https://doi.org/10.48550/arXiv.2606.29776}{10.48550/arXiv.2606.29776}.}

\leavevmode\vadjust pre{\hypertarget{ref-nist_webbook}{}}%
\CSLLeftMargin{23 }%
\CSLRightInline{National Institute of Standards and Technology, NIST
chemistry WebBook, SRD 69, \url{https://webbook.nist.gov}, (accessed 1
January 2026).}

\leavevmode\vadjust pre{\hypertarget{ref-chemotion2024}{}}%
\CSLLeftMargin{24 }%
\CSLRightInline{Chemotion Repository,
\emph{\href{https://doi.org/10.22000/OGoEQGlsZGElrgst}{RADAR4Chem: FT-IR
spectroscopy data}}, Karlsruhe Institute of Technology, 2024.}

\leavevmode\vadjust pre{\hypertarget{ref-sdbs}{}}%
\CSLLeftMargin{25 }%
\CSLRightInline{National Institute of Advanced Industrial Science and
Technology (AIST), Spectral database for organic compounds (SDBS),
\url{https://sdbs.db.aist.go.jp}, (accessed 1 January 2026).}

\leavevmode\vadjust pre{\hypertarget{ref-swain2016chemdataextractor}{}}%
\CSLLeftMargin{26 }%
\CSLRightInline{\href{https://doi.org/10.1021/acs.jcim.6b00207}{M. C.
Swain and J. M. Cole, \emph{J. Chem. Inf. Model.}, 2016, \textbf{56},
1894--1904}.}

\leavevmode\vadjust pre{\hypertarget{ref-pmc_oa}{}}%
\CSLLeftMargin{27 }%
\CSLRightInline{National Library of Medicine, NCBI PMC open access
subset, \url{https://www.ncbi.nlm.nih.gov/pmc/tools/openftlist/},
(accessed 1 January 2026).}

\leavevmode\vadjust pre{\hypertarget{ref-lowe2011opsin}{}}%
\CSLLeftMargin{28 }%
\CSLRightInline{\href{https://doi.org/10.1021/ci100384d}{D. M. Lowe, P.
T. Corbett, P. Murray-Rust and R. C. Glen, \emph{J. Chem. Inf. Model.},
2011, \textbf{51}, 739--753}.}

\leavevmode\vadjust pre{\hypertarget{ref-landrum_rdkit}{}}%
\CSLLeftMargin{29 }%
\CSLRightInline{G. Landrum and others, RDKit
\url{https://www.rdkit.org}.}

\leavevmode\vadjust pre{\hypertarget{ref-krenn2020selfies}{}}%
\CSLLeftMargin{30 }%
\CSLRightInline{\href{https://doi.org/10.1088/2632-2153/aba947}{M.
Krenn, F. Häse, A. Nigam, P. Friederich and A. Aspuru-Guzik, \emph{Mach.
Learn.: Sci. Technol.}, 2020, \textbf{1}, 045024}.}

\leavevmode\vadjust pre{\hypertarget{ref-kim2023pubchem}{}}%
\CSLLeftMargin{31 }%
\CSLRightInline{\href{https://doi.org/10.1093/nar/gkac956}{S. Kim, J.
Chen, T. Cheng and others, \emph{Nucleic Acids Res.}, 2023, \textbf{51},
D1373--D1380}.}

\leavevmode\vadjust pre{\hypertarget{ref-rogers2010ecfp}{}}%
\CSLLeftMargin{32 }%
\CSLRightInline{\href{https://doi.org/10.1021/ci100050t}{D. Rogers and
M. Hahn, \emph{J. Chem. Inf. Model.}, 2010, \textbf{50}, 742--754}.}

\leavevmode\vadjust pre{\hypertarget{ref-mcnemar1947}{}}%
\CSLLeftMargin{33 }%
\CSLRightInline{\href{https://doi.org/10.1007/BF02295996}{Q. McNemar,
\emph{Psychometrika}, 1947, \textbf{12}, 153--157}.}

\leavevmode\vadjust pre{\hypertarget{ref-holm1979}{}}%
\CSLLeftMargin{34 }%
\CSLRightInline{S. Holm, \emph{Scand. J. Stat.}, 1979, \textbf{6},
65--70.}

\leavevmode\vadjust pre{\hypertarget{ref-xu2024contamination}{}}%
\CSLLeftMargin{35 }%
\CSLRightInline{C. Xu, S. Guan, D. Greene and M.-T. Kechadi,
\emph{arXiv}, 2024, preprint, DOI:
\href{https://doi.org/10.48550/arXiv.2406.04244}{10.48550/arXiv.2406.04244}.}

\leavevmode\vadjust pre{\hypertarget{ref-williamson2024mpnn}{}}%
\CSLLeftMargin{36 }%
\CSLRightInline{\href{https://doi.org/10.1016/j.jmr.2024.107795}{D.
Williamson, S. Ponte, I. Iglesias, N. Tonge, C. Cobas and E. K. Kemsley,
\emph{J. Magn. Reson.}, 2024, \textbf{368}, 107795}.}

\leavevmode\vadjust pre{\hypertarget{ref-han2024dftgnn}{}}%
\CSLLeftMargin{37 }%
\CSLRightInline{\href{https://doi.org/10.1021/acs.jctc.4c00422}{C. Han,
D. Zhang, S. Xia and Y. Zhang, \emph{J. Chem. Theory Comput.}, 2024,
\textbf{20}, 5250--5258}.}

\leavevmode\vadjust pre{\hypertarget{ref-xu2025nmrbench}{}}%
\CSLLeftMargin{38 }%
\CSLRightInline{\href{https://doi.org/10.1038/s43588-025-00783-z}{F. Xu,
W. Guo, F. Wang, L. Yao, H. Wang, F. Tang, Z. Gao, L. Zhang, W. E, Z.-Q.
Tian and J. Cheng, \emph{Nat. Comput. Sci.}, 2025, \textbf{5},
292--300}.}

\leavevmode\vadjust pre{\hypertarget{ref-cohen2023delta50}{}}%
\CSLLeftMargin{39 }%
\CSLRightInline{\href{https://doi.org/10.3390/molecules28062449}{R. D.
Cohen, J. S. Wood, Y.-H. Lam, A. V. Buevich, E. C. Sherer, M. Reibarkh,
R. T. Williamson and G. E. Martin, \emph{Molecules}, 2023, \textbf{28},
2449}.}

\leavevmode\vadjust pre{\hypertarget{ref-wang2023selfconsistency}{}}%
\CSLLeftMargin{40 }%
\CSLRightInline{X. Wang, J. Wei, D. Schuurmans, Q. Le, E. Chi, S.
Narang, A. Chowdhery and D. Zhou, in \emph{International conference on
learning representations (ICLR)}, 2023.}

\leavevmode\vadjust pre{\hypertarget{ref-bremser1978hose}{}}%
\CSLLeftMargin{41 }%
\CSLRightInline{\href{https://doi.org/10.1016/S0003-2670(01)83100-7}{W.
Bremser, \emph{Anal. Chim. Acta}, 1978, \textbf{103}, 355--365}.}

\leavevmode\vadjust pre{\hypertarget{ref-kuhn2015nmrshiftdb2}{}}%
\CSLLeftMargin{42 }%
\CSLRightInline{\href{https://doi.org/10.1002/mrc.4263}{S. Kuhn and N.
E. Schlörer, \emph{Magn. Reson. Chem.}, 2015, \textbf{53}, 582--589}.}

\leavevmode\vadjust pre{\hypertarget{ref-praski2025embeddings}{}}%
\CSLLeftMargin{43 }%
\CSLRightInline{M. Praski, J. Adamczyk and W. Czech, \emph{arXiv}, 2025,
preprint, DOI:
\href{https://doi.org/10.48550/arXiv.2508.06199}{10.48550/arXiv.2508.06199}.}

\end{CSLReferences}

\clearpage

\hypertarget{table-of-contents-entry}{%
\section{Table of contents entry}\label{table-of-contents-entry}}

\includegraphics[width=6.3in,height=\textheight]{docs/figures/graphical_abstract.png}

Recall, not verification, is the wall. On real, blind IR + ¹H + ¹³C
literature spectra a frontier LLM proposes the true structure for 34\%
of 194 compounds and, once proposed, selects it 89\% of the time ---
lifting top-1 from 28\% to 30\%.

\end{document}
